\chapter{收敛不只是序列：有限观测的拓扑语言}\label{ch:01}

分析中的对象常由许多局部问题共同描述。一个函数在有限多个点上的值、一个分割对有限多个旧分割的加细、
一个候选测度对有限多个测试函数的回答，都只是有限信息。对每一批有限条件找到一个候选通常并不困难；
困难在于，这些候选是否逼近同一个全局对象，以及极限是否仍然具有连续性、约束条件或其他所需结构。

序列首先遇到一个语法上的障碍。若条件本身没有自然编号，那么“第 (n) 步以后”并不能表达“任意给定的
有限批条件最终同时满足”。有向集、网与滤子将把这句话写成精确的数学陈述。随后出现第二个障碍：
即使每个有限坐标都已有候选答案，也未必有全局极限。紧性和乘积紧性负责产生候选点，半连续性负责在
极限处保留单向不等式。第三个障碍发生在函数空间：逐点极限可能不连续，点态收敛也可能离一致收敛很远。
等度连续性和函数代数分别把有限采样传播为一致控制，并把分离点的观测升级为对所有连续函数的逼近。

本章反复使用下面的四步关系：
\begin{equation}\label{eq:ch1-main-chain}
  \text{有限观测相容}
  \longrightarrow \text{候选极限存在}
  \longrightarrow \text{极限仍属目标类}
  \longrightarrow \text{观测足以识别极限}.
\end{equation}
最后一节讨论可数性、完备性与选择原则。它说明何时可以把网重新压缩为序列，何时需要极大化原理，
以及为什么标准 Borel 空间在概率论中提供了恰当的可数骨架。

本章以本科分析、线性代数、基础点集拓扑以及完备赋范空间的初级结论为基础。复数在本章只作为二维实向量
和标量出现；所需拓扑工具均在本章按使用顺序建立。

为说明某些存在性证明使用了怎样的集合论原则，我们会在证明旁标记 \textsf{ZF}、\textsf{CC}、
\textsf{DC}、\textsf{UFL} 与 \textsf{AC}。这些缩写在 \S\ref{subsec:1-4-3} 正式解释；此前它们只说明
所给证明采用的原则，不增加任何数学假设，也不声称所标强度是最优的。

\section{有向观测：从序列到网与滤子}\label{sec:01-01}

\subsection{有向集、最终性与网}\label{subsec:1-1-1}

序列把逼近过程安放在整数时钟上。这个安排适合逐项部分和，却不适合同时加严许多彼此独立的有限条件。
例如，一次观测也许要求函数在有限个点上接近目标；下一次观测也许换了另一批点。真正需要的“后一步”
不是把某个整数加一，而是把两批点合并。分割的共同加细、有限测试函数的并集和无序和的有限部分和都具有
相同结构。

\begin{example}[有限条件的有向结构]\label{ex:1-1-1-1}
设 \(S\) 为集合，以 \(\Fin(S)\) 表示 \(S\) 的有限子集全体，并按包含关系排序。给定
\(F,G\in\Fin(S)\)，\(F\cup G\) 同时包含二者。若 \((v_s)_{s\in S}\) 位于 Hausdorff 拓扑向量空间，令
\[
  s_F=\sum_{s\in F}v_s,\qquad F\in\Fin(S).
\]
这里每个和都是有限和。若这些有限和在 \(F\) 增大时收敛，则极限没有使用 \(S\) 的编号；这正是无条件
求和所需的语法。拓扑向量空间在本例中只用到有限加法的连续性和极限唯一性。
\end{example}

\begin{definition}[有向集]\label{def:1-1-1-1}
非空集合 \(I\) 上的关系 \(\leq\) 若自反且传递，称为\emph{预序}。若进一步对任意
\(\alpha,\beta\in I\)，存在 \(\gamma\in I\) 使
\[
  \alpha\leq\gamma,\qquad \beta\leq\gamma,
\]
则称 \((I,\leq)\) 为\emph{有向集}。有向关系不要求反对称。
\end{definition}

\begin{definition}[共尾子集]\label{def:1-1-1-cofinal}
有向集 \(I\) 的子集 \(J\) 称为\emph{共尾的}，若对每个 \(\alpha\in I\)，存在 \(\beta\in J\)
满足 \(\alpha\leq\beta\)。这个条件只保证 \(J\) 能越过每个固定阈值；它与定义子网时使用的“最终共尾”
不同。
\end{definition}

\begin{definition}[网、最终与频繁]\label{def:1-1-1-2}
从有向集 \(I\) 到集合 \(X\) 的映射 \(\alpha\mapsto x_\alpha\) 称为 \(X\) 中的\emph{网}。
性质 \(P(x_\alpha)\) \emph{最终成立}，是指存在 \(\alpha_0\in I\)，使每个
\(\alpha\geq\alpha_0\) 都满足 \(P(x_\alpha)\)。性质 \(P(x_\alpha)\) \emph{频繁成立}，是指对每个
\(\alpha_0\in I\)，都存在 \(\alpha\geq\alpha_0\) 满足 \(P(x_\alpha)\)。
\end{definition}

\begin{definition}[网收敛]\label{def:1-1-1-3}
设 \(X\) 为拓扑空间。网 \((x_\alpha)\) \emph{收敛}到 \(x\in X\)，若对 \(x\) 的每个邻域 \(U\)，
关系 \(x_\alpha\in U\) 最终成立。记作 \(x_\alpha\to x\)。
\end{definition}

\begin{definition}[网的聚点]\label{def:1-1-1-cluster}
点 \(x\in X\) 称为网 \((x_\alpha)_{\alpha\in I}\) 的\emph{聚点}，若对 \(x\) 的每个邻域 \(U\)，
关系 \(x_\alpha\in U\) 频繁成立。等价地，对每个 \(\alpha_0\in I\) 和每个
\(U\in\mathcal N(x)\)，
\[
  U\cap\{x_\alpha:\alpha\geq\alpha_0\}\ne\varnothing.
\]
\end{definition}

邻域族本身给出一个重要有向集。规定 \(U\preceq V\) 当且仅当 \(U\supset V\)，指标越大便表示邻域越小，
而 \(U\cap V\) 是共同后继。需要注意，邻域 \(U\) 是 \(X\) 的子集，不是 \(X\) 中的点；从闭包构造网时，
我们将把邻域与其中的点同时放入指标。

\begin{remark}[从序列到 Moore--Smith 网]
十九世纪分析把极限写成序列极限，在度量空间中十分成功。Moore--Smith 网所增加的不是另一套数值计算，
而是把邻域缩小、分割加细和有限条件增多放到同一种“最终同时满足”的语法中。下面的闭包与连续性定理说明，
这种语法恰好恢复了一般拓扑所需的全部极限信息。
\end{remark}

\begin{proposition}[最终性质的有限交法则]\label{proposition:1-1-1-1}
若 \(P(x_\alpha)\) 与 \(Q(x_\alpha)\) 都最终成立，则
\(P(x_\alpha)\wedge Q(x_\alpha)\) 最终成立。同一结论对任意有限多个最终性质成立。
\end{proposition}
\begin{proof}
取阈值 \(\alpha_P,\alpha_Q\)，分别保证 \(P,Q\) 在对应阈值以后成立。由有向性，存在 \(\gamma\) 满足
\(\gamma\geq\alpha_P,\alpha_Q\)。若 \(\alpha\geq\gamma\)，传递性给出
\(\alpha\geq\alpha_P,\alpha_Q\)，故 \(P(x_\alpha)\) 与 \(Q(x_\alpha)\) 同时成立。有限情形由归纳法得到。
\end{proof}

\begin{theorem}[网刻画闭包]\label{theorem:1-1-1-2}
设 \(A\subset X\)，\(x\in X\)。则
\[
  x\in\overline A
  \quad\Longleftrightarrow\quad
  \text{存在取值于 \(A\) 的网收敛到 \(x\)}.
\]
\end{theorem}
\begin{proof}
先设 \(x\in\overline A\)。令
\[
  J=\{(U,a):U\in\mathcal N(x),\ a\in A\cap U\}.
\]
闭包条件保证 \(A\cap U\ne\varnothing\)，所以 \(J\ne\varnothing\)。规定
\((U,a)\preceq(V,b)\) 当且仅当 \(V\subset U\)。给定 \((U,a),(V,b)\)，集合 \(U\cap V\)
是 \(x\) 的邻域，故存在 \(c\in A\cap U\cap V\)，于是 \((U\cap V,c)\) 是共同后继。故 \(J\) 有向。

令 \(z_{(U,a)}=a\)。给定 \(x\) 的邻域 \(W\)，取 \(a_0\in A\cap W\)。当
\((U,a)\succeq(W,a_0)\) 时 \(U\subset W\)，从而 \(z_{(U,a)}=a\in W\)。因此 \(z_j\to x\)。
这里指标集包含所有可行的邻域--点对，没有为整个邻域族预先指定选择函数。

反过来，设 \(A\) 中的网 \(x_\alpha\to x\)。每个 \(x\) 的邻域 \(U\) 最终包含 \(x_\alpha\)，因而
\(A\cap U\ne\varnothing\)。这正是 \(x\in\overline A\)。
\end{proof}

\begin{theorem}[连续性的网判别]\label{theorem:1-1-1-3}
映射 \(f:X\to Y\) 在 \(x\in X\) 连续，当且仅当对每个 \(x_\alpha\to x\)，都有
\(f(x_\alpha)\to f(x)\)。
\end{theorem}
\begin{proof}
若 \(f\) 在 \(x\) 连续，给定 \(f(x)\) 的邻域 \(V\)，存在 \(x\) 的邻域 \(U\) 满足
\(f(U)\subset V\)。由 \(x_\alpha\to x\)，关系 \(x_\alpha\in U\) 最终成立，故
\(f(x_\alpha)\in V\) 最终成立。

反过来，若 \(f\) 在 \(x\) 不连续，则存在 \(f(x)\) 的邻域 \(V\)，使每个 \(x\) 的邻域 \(U\)
都含有某个 \(a\in U\) 满足 \(f(a)\notin V\)。令
\[
  J=\{(U,a):U\in\mathcal N(x),\ a\in U,\ f(a)\notin V\},
\]
并仍按邻域缩小排序。与闭包定理的证明相同，\(J\) 非空且有向，网 \(z_{(U,a)}=a\) 收敛到 \(x\)。
但 \(f(z_j)\notin V\) 对每个 \(j\) 成立，故像网不收敛到 \(f(x)\)，与假设矛盾。
\end{proof}

\begin{proposition}[Hausdorff 判别]\label{proposition:1-1-1-4}
拓扑空间 \(X\) 为 Hausdorff，当且仅当每个网至多有一个极限。
\end{proposition}
\begin{proof}
若 \(X\) 为 Hausdorff，设 \(x_\alpha\to x\) 且 \(x_\alpha\to y\)，其中 \(x\ne y\)。取不交邻域
\(U\ni x,V\ni y\)。网最终进入 \(U\)，也最终进入 \(V\)；由有限交法则，它最终进入
\(U\cap V=\varnothing\)，这是不可能的。

若 \(X\) 不是 Hausdorff，则存在 \(x\ne y\)，使任意邻域 \(U\ni x,V\ni y\) 都相交。令
\[
 J=\{(U,V,z):U\in\mathcal N(x),\ V\in\mathcal N(y),\ z\in U\cap V\}.
\]
按两个邻域同时缩小排序。两个指标的共同后继可从
\((U_1\cap U_2)\cap(V_1\cap V_2)\ne\varnothing\) 中取点得到。令网值为第三坐标。给定 \(x\) 的邻域
\(U_0\)，在指标超过任一以 \(U_0\) 为第一邻域坐标的指标后，网值均落在 \(U_0\)；故该网趋于 \(x\)。
同理它趋于 \(y\)。构造使用全部可行三元组，不需要为所有邻域对预选交点。
\end{proof}

\begin{definition}[第一可数]\label{def:1-1-1-first-countable}
若拓扑空间 \(X\) 的每个点 \(x\) 都有可数邻域基，即存在邻域列 \((V_n)\)，使每个 \(x\) 的邻域都包含
某个 \(V_n\)，则称 \(X\) \emph{第一可数}。
\end{definition}

\begin{proposition}[第一可数空间的序列检验]\label{proposition:1-1-1-first-countable-sequential-tests}
设 \(X\) 第一可数。
\begin{enumerate}
  \item \(x\in\overline A\) 当且仅当存在 \(A\) 中序列收敛到 \(x\)；
  \item 对任意拓扑空间 \(Y\)，映射 \(f:X\to Y\) 在 \(x\) 连续，当且仅当它保持所有趋于 \(x\) 的
  序列之极限。
\end{enumerate}
这并不表示 \(X\) 中每个网都有由 \(\N\) 参数化的子网。
\end{proposition}
\begin{proof}
取 \(x\) 的可数邻域基 \((V_n)\)，并令 \(U_n=V_1\cap\cdots\cap V_n\)。则
\(U_{n+1}\subset U_n\)，且 \((U_n)\) 仍是邻域基。

若 \(x\in\overline A\)，对每个 \(n\) 取 \(a_n\in A\cap U_n\)。给定邻域 \(U\ni x\)，存在 \(N\)
使 \(U_N\subset U\)，于是 \(n\ge N\) 时 \(a_n\in U_n\subset U_N\subset U\)。故 \(a_n\to x\)。
逆向由闭包定义得到。这次逐 \(n\) 选点使用可数选择；若给定邻域基同时给定了选点规则，则不另需选择。

连续映射保持序列极限。反过来，若 \(f\) 在 \(x\) 不连续，取 \(f(x)\) 的邻域 \(V\)，使每个 \(U_n\)
都含有 \(a_n\) 且 \(f(a_n)\notin V\)。上段证明给出 \(a_n\to x\)，而 \(f(a_n)\) 不趋于 \(f(x)\)，矛盾。

最后说明结尾的区别。令 \(I=\Fin(\mathcal P(\N))\)，按包含排序。这里先把稍后
\cref{def:1-1-3-1} 才正式命名的条件写全：映射 \(\phi:\N\to I\) 最终共尾，是指对每个
\(F_0\in I\)，存在 \(N\) 使
\[
  n\ge N\quad\Longrightarrow\quad F_0\subset\phi(n).
\]
若这样的 \(\phi\) 存在，则对每个 \(A\subset\N\)，把 \(F_0=\{A\}\) 代入便有
\(A\in\phi(n)\)（\(n\ge N\)），所以
\[
  \mathcal P(\N)\subset C:=\bigcup_{n\ge1}\phi(n).
\]
对不同的 \(A,B\subset\N\)，规定
\[
  A\prec B
  \quad\Longleftrightarrow\quad
  \min(A\mathbin\triangle B)\in B.
\]
最小差异位置恰属于二者之一，所以任意两个不同集合中恰有一个在前。还需核对传递性：若
\(A\prec B\prec C\)，令
\[
 r=\min(A\mathbin\triangle B),\qquad
 s=\min(B\mathbin\triangle C).
\]
若 \(r=s\)，则 \(A\prec B\) 给出 \(r\in B\)，而 \(B\prec C\) 给出 \(s\notin B\)，矛盾。
若 \(r<s\)，则 \(A,C\) 在所有小于 \(r\) 的位置相同，并且 \(r\notin A\)、\(r\in B=C\)，故
\(A\prec C\)；若 \(s<r\)，同理有 \(s\notin A=B\)、\(s\in C\)，仍得 \(A\prec C\)。因此
\(\prec\) 是 \(\mathcal P(\N)\) 上的严格词典全序。将每个有限集 \(\phi(n)\) 依此序写成
\(A_{n,1},\ldots,A_{n,k_n}\)，再用固定双射 \(\N\leftrightarrow\N^2\)（越界位置填入空集），得到
序列 \((E_j)\) 枚举 \(C\)。令
\[
  B=\{j\in\N:j\notin E_j\}.
\]
则对每个 \(j\) 都有 \(B\ne E_j\)，故 \(B\notin C\)，这与
\(B\in\mathcal P(\N)\subset C\) 矛盾。故这个有向集没有可数最终共尾参数映射；即使取值网是常值网，
也不存在按\cref{def:1-1-3-1}由 \(\N\) 参数化的子网。
\end{proof}

\begin{example}[分割细化]\label{ex:1-1-1-2}
区间上的带标记分割 \((P,\xi)\) 可按底层分割的细化排序。给定两个分割，把两组分点取并，再在每个新区间
内取标签，便得到共同加细。Riemann 和的稳定性要求：存在一个阈值分割，使所有后续细化和所有标签给出的和
都落在规定误差内。本例只说明索引方式；Stieltjes 和是否收敛还需要对被积函数和积分子另加条件。
\end{example}

\begin{example}[闭包中没有序列逼近的点]\label{ex:1-1-1-3}
令 \(I=\mathcal P(\N)\)，\(X=\{0,1\}^{I}\)。暂时给 \(X\) 配置“一个基本邻域只限制有限多个坐标”的
拓扑；这正是 \S\ref{subsec:1-2-2} 将正式定义的乘积拓扑。令 \(A\) 为有限支撑的 \(0/1\) 函数全体，
令 \(\mathbf 1\) 为全一函数。每个只限制有限坐标集 \(F\subset I\) 的 \(\mathbf1\) 邻域都含有
\(\mathbf1_F\in A\)，故 \(\mathbf1\in\overline A\)。网
\[
  F\longmapsto\mathbf1_F,\qquad F\in\Fin(I),
\]
收敛到 \(\mathbf1\)。

另一方面，任取 \(A\) 中序列 \((\mathbf1_{F_n})\)。用上一证明中的规范词典序逐项排列每个有限 \(F_n\)，
再用 \(\N^2\leftrightarrow\N\) 把所得双重有限表写成序列 \((E_j)_{j\ge1}\)；未使用的位置填入
\(\varnothing\)。于是每个 \(F_n\) 的每个元素都在 \((E_j)\) 中出现。令
\[
  i=\{j\in\N:j\notin E_j\}\in I=\mathcal P(\N).
\]
对每个 \(j\) 都有 \(i\ne E_j\)，故 \(i\notin\bigcup_nF_n\)。于是对所有 \(n\)，
\[
  \mathbf1_{F_n}(i)=0\ne1=\mathbf1(i).
\]
集合
\[
  W_i=\{x\in X:x(i)=1\}
\]
只限制坐标 \(i\)，因而按本例给出的基本邻域规则是 \(\mathbf1\) 的开邻域。若
\(\mathbf1_{F_n}\to\mathbf1\)，则应存在 \(N\) 使 \(n\ge N\) 时
\(\mathbf1_{F_n}\in W_i\)，即 \(\mathbf1_{F_n}(i)=1\)；这与上式对每个 \(n\) 的等式矛盾。
因此不必预先调用尚未正式定义的坐标投影连续性，也已证明该序列不收敛到 \(\mathbf1\)。
\end{example}

\begin{figure}[htbp]
\centering
\begin{tikzpicture}[node distance=1.35cm and 1.5cm,>=Stealth,every node/.style={align=center}]
 \node[book node] (n) {邻域缩小\\\(U\cap V\)};
 \node[book node,right=of n] (f) {有限集增大\\\(F\cup G\)};
 \node[book node,right=of f] (p) {分割细化\\公共加细};
 \node[book node,below=of f] (d) {有限阈值有共同后继};
 \node[book node,below=of d] (l) {最终同时满足\\网极限};
 \draw[book arrow] (n)--(d); \draw[book arrow] (f)--(d); \draw[book arrow] (p)--(d);
 \draw[book arrow] (d)--(l);
\end{tikzpicture}
\caption{三个有向系统共享同一个“共同加严”机制。}\label{fig:1-1-1-directed-systems}
\end{figure}

有向关系不必反对称；把预序擅自改成偏序只会加入未使用的条件。“频繁”也不能解释为“无穷多次”，因为
一般指标集没有可数的次数概念。网恢复了序列在一般拓扑中的闭包和连续性角色，却不会把只对可数并成立的
结论推广到任意并。下一节将用网讨论乘积紧性，随后再把同一语言用于函数族极限。

\begin{exercise}
证明有限多个最终性质的合取最终成立。再对序列取 \(P_k(n)\Longleftrightarrow n\ge k\)，说明每个
\(P_k\) 最终成立而 \(\bigwedge_kP_k\) 从不成立。
\end{exercise}
\begin{exercise}
在积有向集上写出“两个变量同时趋近”的定义，并给出两个迭代极限都存在而同时极限不存在的函数例。
\end{exercise}
\begin{exercise}
设 \((v_s)_{s\in S}\) 位于 Banach 空间。证明有限和网收敛，当且仅当对每个 \(\varepsilon>0\)，存在
有限 \(F_0\subset S\)，使每个有限 \(E\subset S\setminus F_0\) 都满足
\(\norm{\sum_{s\in E}v_s}<\varepsilon\)。充分性证明应先对尺度 \(2^{-n}\) 取见证并作有限并，得到递增
\(F_n\)，用序列完备性处理 \(\sum_{s\in F_n}v_s\)，再由同一尾估计回证整个网收敛。
\end{exercise}
\begin{exercise}
证明若同一集合上的两个拓扑拥有完全相同的收敛网，则两个拓扑相同。
\end{exercise}

\subsection{滤子：把尾部行为从索引中分离}\label{subsec:1-1-2}

一个网携带两类信息：参数如何排列，以及哪些集合最终包含网值。很多论证只使用后一类信息。滤子把“最终
落入哪些集合”直接记录成集合族；重参数化造成的冗余随之消失。

\begin{remark}
滤子由 Cartan 学派系统化。网适合写出具体近似，滤子适合合并集合信息；从网到尾滤子会忘掉参数化，
而从滤子返回网若采用成对索引，则无需为每个滤子成员预选一个点。
\end{remark}

\begin{example}[重参数化后尾部信息不变]\label{ex:1-1-2-reparam}
令 \(x_n=1/n\)，并在 \(\N^2\) 的逐坐标次序上令 \(y_{(n,m)}=1/n\)。从 \((n_0,m_0)\) 开始，\(y\)
所能取到的值为 \(\{1/n:n\ge n_0\}\)，与 \(x\) 从 \(n_0\) 开始的尾值相同。第二个时钟没有增加任何
最终信息。
\end{example}

\begin{definition}[滤子]\label{def:1-1-2-1}
集合 \(X\) 上的\emph{滤子}是非空集合族 \(\mathcal F\subset2^X\)，满足
\begin{enumerate}
 \item \(\varnothing\notin\mathcal F\)；
 \item \(A,B\in\mathcal F\) 蕴含 \(A\cap B\in\mathcal F\)；
 \item \(A\in\mathcal F\) 且 \(A\subset B\subset X\) 蕴含 \(B\in\mathcal F\)。
\end{enumerate}
若 \(\mathcal G\supset\mathcal F\)，称 \(\mathcal G\) 比 \(\mathcal F\) 更细。
\end{definition}

\begin{definition}[滤子基]\label{def:1-1-2-2}
非空集合族 \(\mathcal B\subset2^X\) 称为\emph{滤子基}，若
\(\varnothing\notin\mathcal B\)，且对任意 \(B_1,B_2\in\mathcal B\)，存在
\(B_3\in\mathcal B\) 使 \(B_3\subset B_1\cap B_2\)。它生成的滤子为
\[
  \langle\mathcal B\rangle
  =\{A\subset X:\text{存在 }B\in\mathcal B\text{ 使 }B\subset A\}.
\]
\end{definition}

\begin{definition}[滤子收敛与聚点]\label{def:1-1-2-3}
滤子 \(\mathcal F\) \emph{收敛}到 \(x\)，若
\(\mathcal N(x)\subset\mathcal F\)。点 \(x\) 是 \(\mathcal F\) 的\emph{聚点}，若每个
\(F\in\mathcal F\) 与每个 \(U\in\mathcal N(x)\) 都相交。
\end{definition}

这里的 $\mathcal N(x)$ 指所有包含 $x$ 的某个开邻域的集合，而不只指开集本身；它确实是真滤子。
它不含空集，两个成员分别包含的开邻域之交仍是 $x$ 的开邻域，并且“包含一个开邻域”在扩大集合时
保持。因此后面可以把 $\mathcal N(x)$ 代入关于滤子的命题。

\begin{proposition}[尾滤子]\label{proposition:1-1-2-1}
网 \((x_\alpha)_{\alpha\in I}\) 的尾值集合
\(T_\alpha=\{x_\beta:\beta\ge\alpha\}\) 构成滤子基。其生成滤子称为该网的\emph{尾滤子}。
网收敛到 \(x\) 当且仅当尾滤子收敛到 \(x\)；\(x\) 是网的聚点当且仅当它是尾滤子的聚点。
\end{proposition}
\begin{proof}
每个 \(T_\alpha\) 非空。给定 \(T_\alpha,T_\beta\)，取共同后继 \(\gamma\)，则
\(T_\gamma\subset T_\alpha\cap T_\beta\)，故这些尾集是滤子基。

网收敛到 \(x\)，等价于每个 \(U\in\mathcal N(x)\) 包含某个尾集 \(T_\alpha\)，又等价于 \(U\)
属于尾滤子。聚点方面，\(U\) 与每个 \(T_\alpha\) 相交，等价于 \(x_\alpha\in U\) 频繁成立；尾滤子的
任一成员包含某个 \(T_\alpha\)，所以这又等价于 \(U\) 与每个尾滤子成员相交。
\end{proof}

\begin{theorem}[滤子可由网实现]\label{theorem:1-1-2-2}
对 \(X\) 上任一滤子 \(\mathcal F\)，存在 \(X\) 中的网，其尾滤子恰等于 \(\mathcal F\)。构造不需要
为每个滤子成员预选代表点。
\end{theorem}
\begin{proof}
令
\[
 I=\{(F,x):F\in\mathcal F,\ x\in F\},
 \qquad (F,x)\preceq(G,y)\Longleftrightarrow G\subset F.
\]
滤子不含空集，故 \(I\ne\varnothing\)。给定 \((F,x),(G,y)\)，集合
\(F\cap G\in\mathcal F\) 且非空；取 \(z\in F\cap G\)，则 \((F\cap G,z)\) 是共同后继。令
\(z_{(F,x)}=x\)。从任一指标 \((F,x)\) 开始，尾网值全在 \(F\) 中；反过来，每个 \(y\in F\) 都是指标
\((F,y)\) 的网值，且 \((F,x)\preceq(F,y)\)。因此
\[
  \{z_{(G,y)}:(G,y)\succeq(F,x)\}=F.
\]
尾值集合恰遍历 \(\mathcal F\)，故它们生成的向上闭包为
\[
  \{A\subset X:\exists F\in\mathcal F,\ F\subset A\}=\mathcal F.
\]
\end{proof}

\begin{definition}[像滤子]\label{def:1-1-2-image}
若 \(f:X\to Y\)，\(\mathcal F\) 是 \(X\) 上滤子，定义
\[
 f_*\mathcal F=\{B\subset Y:f^{-1}(B)\in\mathcal F\}.
\]
\end{definition}

\begin{proposition}[像滤子的基本性质]\label{proposition:1-1-2-image-properties}
像滤子确为 \(Y\) 上滤子，并且
\[
 f_*\mathcal F
 =\{B\subset Y:\text{存在 }F\in\mathcal F\text{ 使 }f(F)\subset B\}.
\]
若 \(g:Y\to Z\)，则 \((g\circ f)_*=g_*\circ f_*\)。
\end{proposition}
\begin{proof}
因为 \(f^{-1}(\varnothing)=\varnothing\notin\mathcal F\)，像滤子不含空集。逆像保持有限交与包含关系，
所以另两条滤子公理也成立。若 \(f^{-1}(B)\in\mathcal F\)，取 \(F=f^{-1}(B)\)，则 \(f(F)\subset B\)。
反向若 \(F\in\mathcal F\) 且 \(f(F)\subset B\)，则 \(F\subset f^{-1}(B)\)，向上封闭性给
\(f^{-1}(B)\in\mathcal F\)。最后
\((g\circ f)^{-1}(C)=f^{-1}(g^{-1}(C))\)，逐个 \(C\subset Z\) 比较定义即可得到函子性。
\end{proof}

\begin{proposition}[连续性的滤子判别]\label{proposition:1-1-2-3}
\(f:X\to Y\) 在 \(x\) 连续，当且仅当每个收敛到 \(x\) 的滤子 \(\mathcal F\) 都满足
\(f_*\mathcal F\to f(x)\)。
\end{proposition}
\begin{proof}
若 \(f\) 在 \(x\) 连续，给定 \(f(x)\) 的邻域 \(V\)，逆像 \(f^{-1}(V)\) 是 \(x\) 的邻域，故属于
\(\mathcal F\)。于是 \(V\in f_*\mathcal F\)。反过来只取 \(\mathcal F=\mathcal N(x)\)。若
\(f_*\mathcal N(x)\to f(x)\)，则对每个 \(f(x)\) 的邻域 \(V\)，有
\(f^{-1}(V)\in\mathcal N(x)\)，即 \(f^{-1}(V)\) 包含 \(x\) 的某个开邻域。这正是 \(f\) 在 \(x\)
连续。
\end{proof}

\begin{proposition}[聚点与细化]\label{proposition:1-1-2-4}
点 \(x\) 是滤子 \(\mathcal F\) 的聚点，当且仅当存在细化滤子
\(\mathcal G\supset\mathcal F\) 收敛到 \(x\)。
\end{proposition}
\begin{proof}
若 \(x\) 是聚点，集合族
\[
 \mathcal B=\{F\cap U:F\in\mathcal F,\ U\in\mathcal N(x)\}
\]
不含空集。两个基元素的交包含基元素
\((F_1\cap F_2)\cap(U_1\cap U_2)\)，故 \(\mathcal B\) 是滤子基。它生成的滤子 \(\mathcal G\)
同时包含 \(\mathcal F\) 和 \(\mathcal N(x)\)，所以 \(\mathcal G\to x\)。

反过来，若 \(\mathcal G\supset\mathcal F\) 且 \(\mathcal G\to x\)，则对 \(F\in\mathcal F\) 与
\(U\in\mathcal N(x)\)，有 \(F,U\in\mathcal G\)，所以 \(F\cap U\in\mathcal G\)。真滤子不含空集，
故 \(F\cap U\ne\varnothing\)，于是 \(x\) 是 \(\mathcal F\) 的聚点。
\end{proof}

\begin{example}[余有限滤子]\label{ex:1-1-2-1}
无限集合 \(X\) 上所有余有限集构成滤子。对通常序列 \(n\mapsto x_n\)，其尾滤子等于 \(\N\) 上余有限
滤子的像滤子；所以“\(n\to\infty\)”可以理解为一种大集合规则。
\end{example}

\begin{example}[邻域滤子]\label{ex:1-1-2-2}
\(\mathcal N(x)\) 是所有包含某个含 \(x\) 开集的子集所成集合族，而不只是开邻域本身。这样它才向上封闭。
上文已经逐项核对了滤子公理；连续性可写为
\(f_*\mathcal N(x)\supset\mathcal N(f(x))\)。
\end{example}

\begin{example}[概率一事件与依概率收敛]\label{ex:1-1-2-3}
本例只作第~5~章术语预告：概率空间、随机变量、依概率收敛及下面用到的概率运算都将在那里定义和证明，
本章后续不使用本例的结论。
在概率空间上，可测的概率一事件在可测集代数中构成滤子；把它们在 \(\mathcal P(\Omega)\) 中向上封闭，
可得到点集滤子。这表达“忽略零概率例外”的规则。若随机变量 \(X_n\) 依概率收敛到 \(X\)，则对每个
\(\varepsilon>0\)，数值 \(\Prob(\abs{X_n-X}<\varepsilon)\) 趋于 \(1\)；随 \(n\) 变化的这些事件本身
不因此组成滤子。两者的区别在于：概率一事件滤子固定在样本空间上，而依概率收敛控制的是一列事件的概率。
\end{example}

\begin{figure}[htbp]
\centering
\begin{tikzpicture}[node distance=1.45cm and 1.8cm,>=Stealth,every node/.style={align=center}]
 \node[book node] (n) {网\\带参数};
 \node[book node,right=of n] (f) {尾滤子\\最终大集合};
 \node[book node,right=of f] (c) {收敛与聚点\\邻域检验};
 \node[book node,below=of f] (p) {\((F,x)\) 成对索引\\精确恢复};
 \draw[book arrow] (n)--node[above]{忘掉参数}(f);
 \draw[book arrow] (f)--(c);
 \draw[book arrow] (f)--(p);
 \draw[book arrow] (p)--node[below]{\textsf{ZF}}(n);
\end{tikzpicture}
\caption{网与滤子保留相同的最终行为，但滤子忘掉参数化。}\label{fig:1-1-2-net-filter-translation}
\end{figure}

滤子的包含方向容易造成误会：成员越多，条件越强。像滤子不能只写成
\(\{f(F):F\in\mathcal F\}\)，因为这个像集族未必向上封闭。另一方面，给定 \(Y\) 上滤子
\(\mathcal G\)，逆像族 \(\{f^{-1}(G):G\in\mathcal G\}\) 可能含空集；它成为滤子基的充要条件是每个
\(G\) 都与 \(f(X)\) 相交。

\begin{exercise}
证明主滤子 \(\{A\subset X:E\subset A\}\) 为真滤子当且仅当 \(E\ne\varnothing\)。
\end{exercise}
\begin{exercise}
证明逆像族 \(\{f^{-1}(G):G\in\mathcal G\}\) 为滤子基当且仅当每个 \(G\in\mathcal G\) 与
\(f(X)\) 相交，并用常值映射构造失败例。
\end{exercise}
\begin{exercise}
证明滤子在 Hausdorff 空间中至多收敛到一点。
\end{exercise}
\begin{exercise}
用 \((F,x)\) 成对索引从给定滤子构造尾滤子恰等于它的网，并与逐 \(F\) 预选 \(x_F\in F\) 的写法比较。
\end{exercise}

滤子已经把“最终落入”改写为集合族的包含关系；要从一个不收敛的对象中提取稳定部分，还需说明何种
细化既保留全部尾部信息，又足以在每个集合及其补集之间作出选择。

\subsection{子网、超滤子与选择原则}\label{subsec:1-1-3}

从近似族中提取稳定部分，是分析中反复出现的动作。序列使用子序列；一般网需要一种既能越过每个旧阈值、
又不会不断返回早期指标的重参数化。滤子的对应动作是细化。若进一步要求每个集合 \(A\) 与其补集之间必须
作出一致选择，就得到超滤子。

\begin{remark}
子网在文献中有若干等价或近似的定义。这里采用最终共尾映射，因为它保证最终性质传递，并保证聚点能够产生
收敛子网。超滤子引理弱于完整选择公理这一集合论事实将在 \S\ref{subsec:1-4-3} 中准确表述；本小节只在
需要扩张滤子时明确说明所采用的原则。
\end{remark}

\begin{definition}[子网]\label{def:1-1-3-1}
网 \((y_\beta)_{\beta\in J}\) 称为 \((x_\alpha)_{\alpha\in I}\) 的\emph{子网}，若存在映射
\(\phi:J\to I\)，使 \(y_\beta=x_{\phi(\beta)}\)，且 \(\phi\) \emph{最终共尾}：对每个
\(\alpha_0\in I\)，存在 \(\beta_0\in J\)，使
\[
 \beta\ge\beta_0\quad\Longrightarrow\quad\phi(\beta)\ge\alpha_0.
\]
\end{definition}

\begin{example}[仅有共尾像还不够]\label{ex:1-1-3-3}
令 \(x_n=1/n\)，并定义 \(\phi(2k)=k\)、\(\phi(2k-1)=1\)。像集 \(\phi(\N)\) 在 \(\N\) 中共尾，
但 \(\phi\) 不是最终共尾：每个奇数指标仍映到 \(1\)。重参数后的网在奇数项恒取 \(1\)，因而不趋于 \(0\)。
\end{example}

\begin{proposition}[子网继承最终性质]\label{proposition:1-1-3-subnet-eventual}
若 \(P(x_\alpha)\) 最终成立，且 \((x_{\phi(\beta)})\) 是子网，则
\(P(x_{\phi(\beta)})\) 最终成立。特别地，收敛网的每个子网收敛到同一给定极限。
\end{proposition}
\begin{proof}
取 \(\alpha_0\) 使 \(\alpha\ge\alpha_0\) 时 \(P(x_\alpha)\) 成立。最终共尾性给出 \(\beta_0\)，使
\(\beta\ge\beta_0\) 时 \(\phi(\beta)\ge\alpha_0\)。于是这些 \(\beta\) 都满足
\(P(x_{\phi(\beta)})\)。对收敛网，把 \(P\) 取为“网值属于给定极限的邻域”即可。
\end{proof}

\begin{definition}[超滤子]\label{def:1-1-3-2}
真滤子 \(\mathcal U\) 若在包含关系下极大，称为\emph{超滤子}。
\end{definition}

\begin{definition}[超滤子极限]\label{def:1-1-3-3}
若超滤子 \(\mathcal U\) 收敛到 \(x\)，称 \(x\) 为其极限。
\end{definition}

\begin{proposition}[超滤子二分刻画]\label{proposition:1-1-3-2}
真滤子 \(\mathcal U\) 为超滤子，当且仅当对每个 \(A\subset X\)，恰有 \(A\in\mathcal U\) 或
\(A^c\in\mathcal U\) 之一成立。
\end{proposition}
\begin{proof}
先设 \(\mathcal U\) 极大。若 \(A\notin\mathcal U\) 且对每个 \(U\in\mathcal U\) 都有
\(A\cap U\ne\varnothing\)，则
\[
  \mathcal B_A=\{A\cap U:U\in\mathcal U\}
\]
不含空集，并且
\[
  (A\cap U_1)\cap(A\cap U_2)=A\cap(U_1\cap U_2)\in\mathcal B_A.
\]
故 \(\mathcal B_A\) 是滤子基；它生成的真滤子同时包含 \(\mathcal U\) 与 \(A\)，从而严格细化
\(\mathcal U\)，与极大性矛盾。因此存在 \(U\in\mathcal U\) 使 \(A\cap U=\varnothing\)。于是
\(U\subset A^c\)，向上封闭性给出 \(A^c\in\mathcal U\)。二者不能同时属于 \(\mathcal U\)，因为
\[
  A,A^c\in\mathcal U\quad\Longrightarrow\quad
  \varnothing=A\cap A^c\in\mathcal U.
\]

反过来，设二分性质成立，且真滤子 \(\mathcal G\supset\mathcal U\)。若存在
\(A\in\mathcal G\setminus\mathcal U\)，则 \(A^c\in\mathcal U\subset\mathcal G\)，于是
\(\varnothing=A\cap A^c\in\mathcal G\)，矛盾。因此 \(\mathcal G=\mathcal U\)，即 \(\mathcal U\) 极大。
\end{proof}

\begin{remark}
超滤子的聚点自动成为极限。若 \(x\) 是 \(\mathcal U\) 的聚点而 \(U\in\mathcal N(x)\) 不属于
\(\mathcal U\)，二分性质给 \(U^c\in\mathcal U\)，但 \(U^c\cap U=\varnothing\)，与聚点定义矛盾。
\end{remark}

\begin{example}[自由超滤子]\label{ex:1-1-3-1}
无限集合 \(X\) 上的余有限滤子若扩张为超滤子 \(\mathcal U\)，则 \(\mathcal U\) 不含有限集。事实上，
若有限集 \(F\in\mathcal U\)，则余有限集 \(X\setminus F\) 也属于 \(\mathcal U\)，从而
\[
  \varnothing=F\cap(X\setminus F)\in\mathcal U,
\]
矛盾。这样的 \(\mathcal U\) 称为自由超滤子。它的存在来自超滤子引理；这里不把某个自由超滤子当作
可显式列出的对象。
\end{example}

\begin{example}[有限划分与可数划分]\label{ex:1-1-3-2}
任一 \(\N\) 上超滤子恰包含偶数集或奇数集之一。更一般地，若
\(X=A_1\sqcup\cdots\sqcup A_m\)，且没有 \(A_j\) 属于 \(\mathcal U\)，二分性质便给出
\(A_j^c\in\mathcal U\)（\(1\le j\le m\)），从而
\[
  \varnothing=\bigcap_{j=1}^m A_j^c\in\mathcal U,
\]
矛盾；而两个不同胞腔不能同时属于 \(\mathcal U\)，因为它们的交为空。因此有限划分中恰有一个胞腔
属于超滤子。有限性不能删除：自由超滤子包含 \(\N\)，却不包含任何单点，因而不选择可数单点划分中的
胞腔。
\end{example}

\begin{theorem}[聚点--子网定理]\label{theorem:1-1-3-1}
点 \(x\) 是网 \((x_\alpha)_{\alpha\in I}\) 的聚点，当且仅当该网有一个收敛到 \(x\) 的子网。
\end{theorem}
\begin{proof}
若有子网 \(x_{\phi(\beta)}\to x\)，给定 \(x\) 的邻域 \(U\) 和旧阈值 \(\alpha_0\)，子网收敛给出
\(\beta_1\)，使 \(\beta\ge\beta_1\) 时 \(x_{\phi(\beta)}\in U\)；最终共尾性给出 \(\beta_2\)，使
\(\beta\ge\beta_2\) 时 \(\phi(\beta)\ge\alpha_0\)。取两者的共同后继 \(\beta\)，便得到一个
\(\alpha=\phi(\beta)\ge\alpha_0\) 且 \(x_\alpha\in U\)。故 \(x\) 是聚点。

反过来设 \(x\) 是聚点。令
\[
 J=\{(\alpha,U):\alpha\in I,\ U\in\mathcal N(x),\ x_\alpha\in U\},
\]
并定义
\[
 (\alpha,U)\preceq(\beta,V)
 \quad\Longleftrightarrow\quad
 \alpha\le\beta\text{ 且 }V\subset U.
\]
聚点定义保证 \(J\ne\varnothing\)。给定 \((\alpha,U),(\beta,V)\)，先取旧指标共同后继 \(\delta\)。
因为网频繁进入 \(U\cap V\)，存在 \(\gamma\ge\delta\) 使 \(x_\gamma\in U\cap V\)。于是
\((\gamma,U\cap V)\) 是共同后继，故 \(J\) 有向。

令 \(\phi(\alpha,U)=\alpha\)。给定 \(\alpha_0\)，取任意邻域 \(W\)；聚点性给出
\(\alpha_1\ge\alpha_0\) 且 \(x_{\alpha_1}\in W\)，所以 \((\alpha_1,W)\in J\)。若
\((\beta,V)\succeq(\alpha_1,W)\)，则 \(\phi(\beta,V)=\beta\ge\alpha_1\ge\alpha_0\)。故
\(\phi\) 最终共尾。

子网值为 \(x_{\phi(\alpha,U)}=x_\alpha\)。给定邻域 \(W\ni x\)，取 \((\alpha_1,W)\in J\) 如上；
在该指标之后，邻域坐标 \(V\subset W\)，且子网值 \(x_\beta\in V\subset W\)。故子网收敛到 \(x\)。
整个构造使用所有可行对，属于 \textsf{ZF} 证明。
\end{proof}

下面的证明使用一条外部集合论原则。Zorn 引理断言：若一个非空偏序集的每条链都有上界，则该偏序集
有极大元。在 \textsf{ZF} 上，Zorn 引理与选择公理 \textsf{AC} 等价；本书不证明这项集合论等价，
只在下述证明中明确调用 Zorn 引理。超滤子引理本身可以作为较弱的 \textsf{UFL} 原则单独采用。

\begin{theorem}[超滤子引理]\label{theorem:1-1-3-3}
每个真滤子都包含于某个超滤子。
\end{theorem}
\begin{proof}
在通常采用选择公理的语境中，可由 Zorn 引理证明。固定真滤子 \(\mathcal F\)，令 \(\mathscr P\) 为所有
包含 \(\mathcal F\) 的真滤子，按包含排序。若 \(\mathscr C\subset\mathscr P\) 是链，令
\(\mathcal H=\bigcup_{\mathcal G\in\mathscr C}\mathcal G\)。空集不属于 \(\mathcal H\)。若
\(A,B\in\mathcal H\)，存在链中两个滤子分别包含它们；链的线性次序使其中一个滤子同时包含 \(A,B\)，
所以 \(A\cap B\in\mathcal H\)。向上封闭性也由同一成员滤子继承。因此 \(\mathcal H\) 是链的上界。
Zorn 引理给出极大成员，它就是包含 \(\mathcal F\) 的超滤子。

这份证明使用 \textsf{AC}。命题本身通常称超滤子引理，并可作为较弱原则 \textsf{UFL} 采用；上述证明
不说明命题必须使用完整选择公理。
\end{proof}

\begin{proposition}[超滤子像]\label{proposition:1-1-3-4}
若 \(f:X\to Y\)，\(\mathcal U\) 是 \(X\) 上超滤子，则 \(f_*\mathcal U\) 是 \(Y\) 上超滤子。
\end{proposition}
\begin{proof}
像滤子是真滤子。对任意 \(B\subset Y\)，超滤子的二分性质应用于 \(f^{-1}(B)\) 给出
\(f^{-1}(B)\in\mathcal U\) 或
\(X\setminus f^{-1}(B)=f^{-1}(B^c)\in\mathcal U\)，且二者不能同时成立。依像滤子定义，
\(B,B^c\) 中恰有一个属于 \(f_*\mathcal U\)。由二分刻画，像滤子为超滤子。
\end{proof}

\begin{figure}[htbp]
\centering
\begin{tikzpicture}[node distance=1.45cm and 1.8cm,>=Stealth,every node/.style={align=center}]
 \node[book node] (a) {网的聚点};
 \node[book node,right=of a] (b) {收敛子网};
 \node[book node,below=of a] (c) {尾滤子的聚点};
 \node[book node,right=of c] (d) {收敛细化};
 \node[book node,below right=of c] (e) {超滤子扩张};
 \draw[book arrow,<->] (a)--node[above]{\textsf{ZF}}(b);
 \draw[book arrow,<->] (a)--node[left]{翻译}(c);
 \draw[book arrow,<->] (c)--node[above]{\textsf{ZF}}(d);
 \draw[book dashed arrow] (d)--node[right]{\textsf{UFL}}(e);
\end{tikzpicture}
\caption{稳定部分的三种语言以及极大化发生的位置。}\label{fig:1-1-3-cluster-subnet-ultrafilter}
\end{figure}

子网不必由原指标集的子集得到；关键是最终共尾映射。超滤子极限也不自动唯一，唯一性仍需 Hausdorff 条件。
超滤子的二分只控制有限划分，这一点将在测度论的可数可加性出现时形成重要对照。

\begin{exercise}
证明收敛网的每个子网收敛到同一给定极限。
\end{exercise}
\begin{exercise}
证明对超滤子 \(\mathcal U\)，有限并 \(A_1\cup\cdots\cup A_n\in\mathcal U\) 当且仅当某个
\(A_j\in\mathcal U\)。
\end{exercise}
\begin{exercise}
从网的聚点出发，重新构造聚点--子网定理中的有向集，并逐项检查有向性、最终共尾性与收敛性。
\end{exercise}
\begin{exercise}
设 \(X\) 第一可数。对普通序列证明：\(x\) 是序列聚点，当且仅当存在收敛到 \(x\) 的子序列。
\end{exercise}

下一节把紧性表述为：所有有向观测都能产生聚点，或者所有超滤子都能找到极限。

\section{紧性、乘积与极小化}\label{sec:01-02}

\subsection{紧性：有限交、网与超滤子的统一}\label{subsec:1-2-1}

在欧氏空间里，Bolzano--Weierstrass 定理很容易使人把紧性理解为“每个序列都有收敛子列”。一般拓扑中，
序列只观察可数长的路线，而开覆盖可能同时包含不可数多个局部条件。下面的乘积首先暴露这个缺口。

\begin{example}[一个没有收敛子列的坐标序列]\label{ex:1-2-1-1}
令
\[
 X=\{0,1\}^{\mathcal P(\N)},\qquad x_n(A)=\1_A(n).
\]
任取子序列 \((x_{n_k})\)。令 \(A=\{n_{2k}:k\ge1\}\)，则
\[
 x_{n_k}(A)=
 \begin{cases}
  0,&k\text{ 为奇数},\\
  1,&k\text{ 为偶数}.
 \end{cases}
\]
因此该子序列在第 \(A\) 个坐标不收敛，也就不可能逐坐标收敛。到
\S\ref{subsec:1-2-2} 证明 Tychonoff 定理以后，我们会知道 \(X\) 在乘积拓扑下紧；在此之前只使用上面的
非收敛计算。
\end{example}

\begin{definition}[紧空间]\label{def:1-2-1-1}
拓扑空间 \(X\) 称为\emph{紧的}，若每个开覆盖都有有限子覆盖。闭集族具有\emph{有限交性质}，是指
其每个有限子族（包括空子族）的交非空；空子族的交约定为 \(X\)。
\end{definition}

\begin{proposition}[开覆盖与有限交性质]\label{proposition:1-2-1-cover-fip}
空间 \(X\) 紧，当且仅当每个具有有限交性质的闭集族都有非空交。这个约定包括空空间的情形。
\end{proposition}
\begin{proof}
设 \(X\) 紧，闭集族 \((F_i)_{i\in I}\) 的总交为空。则 \((X\setminus F_i)_{i\in I}\) 是开覆盖，
故有有限子覆盖。取补集可知相应有限多个 \(F_i\) 的交为空，所以该闭集族不具有有限交性质。

反过来，若开覆盖 \((U_i)_{i\in I}\) 没有有限子覆盖，则每个有限 \(E\subset I\) 都满足
\[
 \bigcap_{i\in E}(X\setminus U_i)
 =X\setminus\bigcup_{i\in E}U_i\ne\varnothing.
\]
于是闭集族 \((X\setminus U_i)\) 具有有限交性质。若假设中的闭集结论成立，其总交非空；但总交是
\(X\setminus\bigcup_iU_i=\varnothing\)，矛盾。
\end{proof}

有限交性质与滤子基不是同一个条件。前者只要求原集合族的每个有限交非空；后者还要求对任意两个基元素，
原集合族中已有另一个元素包含在二者之交内。例如在 \(X=\{1,2,3\}\) 上，
\[
  \mathcal C=\bigl\{\{1,2\},\{2,3\}\bigr\}
\]
具有有限交性质，却不是滤子基，因为 \(\mathcal C\) 中没有元素包含于 \(\{2\}\)。若 \(\mathcal C\)
具有有限交性质，则把它的所有有限交加入后，所得集合族便构成滤子基，并生成一个真滤子。这个加闭步骤正是
后面把闭集有限交性质转成滤子或超滤子时所用的操作。

\begin{definition}[相对紧]\label{def:1-2-1-relative-compact}
对子集 \(A\subset X\)，若其在环境空间 \(X\) 中的闭包 \(\overline A^{\,X}\) 紧，则称 \(A\) 在 \(X\)
中\emph{相对紧}。环境空间是定义的一部分；例如 \((0,1)\) 在 \(\R\) 中相对紧，而把 \((0,1)\) 自身作为
环境时，它的闭包仍为非紧空间 \((0,1)\)。
\end{definition}

\begin{proposition}[聚点的尾闭包表示]\label{proposition:1-2-1-tail-closure}
对网 \((x_\alpha)_{\alpha\in I}\)，令
\[
 K_\alpha=\overline{\{x_\beta:\beta\ge\alpha\}}.
\]
则该网的聚点集合为 \(\bigcap_{\alpha\in I}K_\alpha\)。
\end{proposition}
\begin{proof}
若 \(x\) 是聚点，给定 \(\alpha\) 与 \(x\) 的邻域 \(U\)，频繁进入性质给出
\(\beta\ge\alpha\) 使 \(x_\beta\in U\)。故 \(x\in K_\alpha\)，对每个 \(\alpha\) 都成立。

反过来，若 \(x\in\bigcap_\alpha K_\alpha\)，则对任意阈值 \(\alpha\) 与邻域 \(U\ni x\)，闭包定义给出
\(U\cap\{x_\beta:\beta\ge\alpha\}\ne\varnothing\)。所以网频繁进入每个 \(x\) 的邻域，\(x\) 是聚点。
\end{proof}

\begin{definition}[序列紧与可数紧]\label{def:1-2-1-3}
空间称为\emph{序列紧的}，若每个序列都有收敛子序列；称为\emph{可数紧的}，若每个可数开覆盖都有有限
子覆盖。一般拓扑空间中，这两个性质均不等同于紧性。
\end{definition}

\begin{theorem}[网紧性定理]\label{theorem:1-2-1-1}
空间 \(X\) 紧，当且仅当每个 \(X\) 中的网都有聚点；等价地，每个网都有收敛子网。
\end{theorem}
\begin{proof}
先设 \(X\) 紧。对网 \((x_\alpha)\) 考虑上一命题的尾闭包 \(K_\alpha\)。给定有限多个指标，取其共同后继
\(\gamma\)，则 \(K_\gamma\) 包含于相应有限交，故 \((K_\alpha)\) 具有有限交性质。紧性给出
\(x\in\bigcap_\alpha K_\alpha\)，上一命题说明 \(x\) 是聚点。由
\cref{theorem:1-1-3-1}，网有收敛子网。

反过来，设每个网都有聚点。令 \((F_i)_{i\in I}\) 是具有有限交性质的闭集族。若这个情形发生，则
\(X\ne\varnothing\)。令
\[
 J=\{(E,x):E\subset I\text{ 有限},\ x\in\bigcap_{i\in E}F_i\},
\]
并按有限集合坐标增大排序。有限交性质保证 \(J\ne\varnothing\)。两个指标的有限集合坐标取并，再从对应
非空交中取一点，得到共同后继，所以 \(J\) 有向。令网值 \(z_{(E,x)}=x\)，并取其聚点 \(z\)。

固定 \(i\in I\)。在任何有限集合坐标含 \(i\) 的指标之后，所有网值都落在 \(F_i\)。若 \(z\notin F_i\)，
则开邻域 \(X\setminus F_i\) 不可能被该尾部访问，与 \(z\) 是聚点矛盾。因此 \(z\in F_i\) 对每个 \(i\)
成立，即 \(\bigcap_iF_i\ne\varnothing\)。由有限交刻画，\(X\) 紧。这个反向构造使用全部可行对，不为每个
有限交预选一点，属于 \textsf{ZF} 证明。

“每个网有收敛子网”蕴含“每个网有聚点”由聚点--子网定理的逆向得到，故三种陈述等价。
\end{proof}

\begin{theorem}[超滤子紧性定理]\label{theorem:1-2-1-2}
在采用超滤子引理的语境中，\(X\) 紧，当且仅当 \(X\) 上每个超滤子至少收敛到一点。
\end{theorem}
\begin{proof}
先设 \(X\) 紧，\(\mathcal U\) 是超滤子。闭集族
\(\{\overline U:U\in\mathcal U\}\) 具有有限交性质，因为有限多个 \(U\) 的交仍属于真滤子且非空。
取
\[
 x\in\bigcap_{U\in\mathcal U}\overline U.
\]
对每个 \(V\in\mathcal N(x)\) 和 \(U\in\mathcal U\)，由 \(x\in\overline U\) 得 \(V\cap U\ne\varnothing\)，
所以 \(x\) 是 \(\mathcal U\) 的聚点。超滤子的聚点是极限，故 \(\mathcal U\to x\)。这一方向不需扩张
滤子，属于 \textsf{ZF} 证明。

反过来，设每个超滤子都有极限，令 \((F_i)\) 为具有有限交性质的闭集族。有限交组成滤子基；由超滤子引理，
其生成滤子可扩张为超滤子 \(\mathcal U\)。取 \(\mathcal U\to x\)。每个 \(F_i\in\mathcal U\)。若
\(x\notin F_i\)，则 \(X\setminus F_i\) 是 \(x\) 的邻域，也属于 \(\mathcal U\)，导致空集属于
\(\mathcal U\)。故 \(x\in\bigcap_iF_i\)，由有限交刻画可知 \(X\) 紧。这一方向使用 \textsf{UFL}。
\end{proof}

\begin{proposition}[紧集在 Hausdorff 空间中闭]\label{proposition:1-2-1-3}
若 \(K\subset X\) 紧且 \(X\) Hausdorff，则 \(K\) 闭。Hausdorff 不能只替换为 \(T_1\)。
\end{proposition}
\begin{proof}
取 \(x\in X\setminus K\)。令 \(\mathscr U\) 为所有满足下列条件的开集 \(U\)：存在 \(x\) 的开邻域
\(V\) 使 \(U\cap V=\varnothing\)。Hausdorff 性保证每个 \(y\in K\) 属于某个 \(U\in\mathscr U\)，
所以 \(\mathscr U\) 覆盖 \(K\)。紧性给出有限子覆盖 \(U_1,\ldots,U_m\)。对这有限多个 \(U_j\)，分别取
\(x\) 的开邻域 \(V_j\) 使 \(U_j\cap V_j=\varnothing\)。则 \(V=\bigcap_jV_j\) 是 \(x\) 的开邻域，
并与 \(\bigcup_jU_j\supset K\) 不交。因此 \(X\setminus K\) 开，\(K\) 闭。

反例取无限集合上的余有限拓扑。它是 \(T_1\) 且紧：任取开覆盖，先取一个非空成员，其补集有限，再为剩余
有限多个点各取一个覆盖成员即可。任一无限真子集在子空间余有限拓扑下也紧，却不是环境空间中的闭集。
\end{proof}

\begin{example}[递减闭集网]\label{ex:1-2-1-2}
若 \(X\) 紧，\((F_\alpha)\) 是按指标递减的非空闭集网，则这些闭集具有有限交性质，所以
\(\bigcap_\alpha F_\alpha\ne\varnothing\)。可算模型是
\(\bigcap_n[0,1/n]=\{0\}\)。若删去紧性，\(\R\) 中 \(F_n=[n,\infty)\) 每个都非空闭且递减，总交却为空：
约束的满足点逃向无穷。
\end{example}

\begin{theorem}[度量空间中的序列判别]\label{theorem:1-2-1-4}
度量空间 \(X\) 紧，当且仅当 \(X\) 序列紧。
\end{theorem}
\begin{proof}
若 \(X=\varnothing\)，则开覆盖紧性与序列紧性都成立。以下设 \(X\ne\varnothing\)。
设 \(X\) 紧，\((x_n)\) 是序列。把它看成网，网紧性定理给出聚点 \(x\)。递归取严格递增整数 \(n_k\)，使
\[
 n_k>n_{k-1},\qquad d(x_{n_k},x)<1/k.
\]
聚点性质保证每一步都有可选的 \(n_k\)。于是 \(x_{n_k}\to x\)。这份递归抽取使用 \textsf{DC}。

反过来设 \(X\) 序列紧。先证全有界。若存在 \(\varepsilon>0\) 使 \(X\) 不能由有限多个
\(\varepsilon\)-球覆盖，可递归选择 \(x_{n+1}\) 不属于
\(\bigcup_{j\le n}B(x_j,\varepsilon)\)。所得序列两两距离至少 \(\varepsilon\)，所以没有 Cauchy 子列，
更没有收敛子列，矛盾。

对每个 \(m\ge1\)，取有限 \(1/m\)-网 \(D_m\)，则
\[
 \mathcal B=\{B(a,r):a\in\bigcup_mD_m,\ r\in\Q,\ r>0\}
\]
是可数基。事实上，若 \(x\in U\) 且 \(U\) 开，取 \(\rho>0\) 使 \(B(x,\rho)\subset U\)，再取
\(m\) 使 \(2/m<\rho\) 及 \(a\in D_m\) 使 \(d(a,x)<1/m\)，并选有理
\(r\) 满足 \(1/m<r<\rho-1/m\)。于是 \(x\in B(a,r)\subset U\)。

给定开覆盖 \(\mathscr U\)，对每个被某个覆盖成员包含的基元素，选取一个这样的覆盖成员。所有选出的成员
构成可数子覆盖；把它编号为 \(U_1,U_2,\ldots\)。若它没有有限子覆盖，逐 \(n\) 取
\[
 x_n\in X\setminus\bigcup_{j=1}^nU_j.
\]
取收敛子列 \(x_{n_k}\to x\)，并选 \(m\) 使 \(x\in U_m\)。当 \(k\) 充分大时
\(x_{n_k}\in U_m\)，而 \(n_k\ge m\) 时定义又给 \(x_{n_k}\notin U_m\)，矛盾。因此原开覆盖有有限
子覆盖。逐尺度取有限网、从可数子覆盖取反例点和递归抽取所使用的原则不超过通常的
\textsf{CC}/\textsf{DC}；这里不主张这是命题的最优集合论强度。
\end{proof}

\begin{example}[紧性、序列紧性与全有界性]\label{ex:1-2-1-3}
在度量空间中，上一证明说明开覆盖紧性和序列紧性重新一致。结合完备性与全有界性还可得到第三种等价刻画；
其证明留作本小节习题，以便把一般拓扑中的区别与度量特例分开。

全有界性不能只换成有界性。例如 \(\ell^2\) 的闭单位球完备且有界，但标准正交向量满足
\(\norm{e_m-e_n}=\sqrt2\)（\(m\ne n\)）。因而半径小于 \(\sqrt2/2\) 的有限多个球不能覆盖该闭单位球；
它不是全有界的，也就不是紧的。
\end{example}

\begin{figure}[htbp]
\centering
\begin{tikzpicture}[node distance=1.35cm and 1.7cm,>=Stealth,every node/.style={align=center}]
 \node[book strong node] (a) {紧性\\有限子覆盖};
 \node[book node,left=of a] (b) {闭集 FIP\\交非空};
 \node[book node,right=of a] (c) {网有聚点\\有收敛子网};
 \node[book node,below=of a] (d) {每个超滤子\\有极限};
 \draw[book arrow,<->] (a)--node[above]{补集}(b);
 \draw[book arrow,<->] (a)--node[above]{\textsf{ZF}}(c);
 \draw[book arrow] (a)--node[left]{\textsf{ZF}}(d);
 \draw[book dashed arrow] (d)--node[right]{反向用 \textsf{UFL}}(a);
\end{tikzpicture}
\caption{紧性的四种语言以及超滤子反向证明所需的扩张原则。}
\label{fig:1-2-1-compactness-equivalences}
\end{figure}

紧性给出超滤子极限的存在，却不保证唯一；唯一性仍由 Hausdorff 性提供。相对紧也始终相对于环境空间定义。
在一般拓扑中，序列紧与可数紧都不能替代紧性；度量结构提供的可数基才使序列重新足够。

\begin{exercise}
证明度量空间紧当且仅当它完备且全有界。两个方向都要写出，并说明从全有界序列抽取 Cauchy 子列时
怎样逐层选择有限网中的嵌套无限子列。
\end{exercise}
\begin{exercise}
不使用 Tychonoff 定理，证明有限个紧空间的乘积紧。
\end{exercise}
\begin{exercise}
证明紧空间到 Hausdorff 空间的连续双射是同胚。
\end{exercise}
\begin{exercise}
令 \(\omega_1\) 为第一不可数序数，并给序数区间 \([0,\omega_1)\) 配置序拓扑。证明每个可数子集
都有严格小于 \(\omega_1\) 的上界，并据此证明该空间既可数紧又序列紧，但不紧。最后证明它是
Hausdorff 空间，并指出论证中何处使用了“可数个可数序数的上确界仍是可数序数”。
本题是集合论拓扑的选读旁例；序数、\(\omega_1\)、序拓扑以及引号内的集合论结论在本题中作为外部
定义与待证引理，本章正文后续不依赖本题。
\end{exercise}
\begin{exercise}
证明在紧 Hausdorff 空间中，网收敛到 \(x\) 当且仅当 \(x\) 是它的唯一聚点。指出“唯一聚点蕴含收敛”
使用紧性，而排除第二聚点使用 Hausdorff 性。
\end{exercise}

这些等价刻画处理的是单个空间。函数族与坐标空间把许多局部紧性问题同时放在各个坐标上，因而接下来要
判断：逐坐标的有限观测能否重新拼成整体紧性。

\subsection{乘积拓扑、逐点收敛与 Tychonoff 定理}\label{subsec:1-2-2}

一个函数可以看作其全部函数值组成的坐标族。若一次实验只能检查有限多个输入点，那么合适的邻域也只应
限制有限多个坐标。允许一次限制所有坐标会得到更细的盒拓扑；下面的失败例将说明，那种拓扑既不描述点态
收敛，也不保留乘积紧性。

\begin{remark}
有限乘积紧性可以反复使用两因子情形得到；任意乘积则把选择原则带入拓扑核心。超滤子证明会把两件事分开：
因子紧性产生坐标极限，而有限柱只负责把坐标收敛拼成乘积收敛。
\end{remark}

\begin{definition}[乘积拓扑]\label{def:1-2-2-1}
对空间族 \((X_i)_{i\in I}\)，乘积 \(X=\prod_{i\in I}X_i\) 上的\emph{乘积拓扑}是使每个坐标投影
\(\pi_i:X\to X_i\) 连续的最弱拓扑。它的基本开集为
\[
 \bigcap_{i\in F}\pi_i^{-1}(U_i),
\]
其中 \(F\subset I\) 有限，且每个 \(U_i\subset X_i\) 开。
\end{definition}

\begin{definition}[逐点收敛拓扑]\label{def:1-2-2-2}
函数空间 \(Y^S=\prod_{s\in S}Y\) 上的乘积拓扑称为\emph{逐点收敛拓扑}。
\end{definition}

\begin{definition}[柱集]\label{def:1-2-2-3}
只由有限多个坐标条件决定的集合称为\emph{柱集}。坐标条件均开时，所得开柱构成乘积拓扑的基。
\end{definition}

\begin{theorem}[逐坐标收敛判别]\label{theorem:1-2-2-1}
网 \(x_\alpha\in\prod_iX_i\) 收敛到 \(x\)，当且仅当对每个 \(i\in I\)，
\(\pi_i(x_\alpha)\to\pi_i(x)\)。
\end{theorem}
\begin{proof}
若 \(x_\alpha\to x\)，坐标投影连续，所以每个坐标网收敛。

反过来，设每个坐标网收敛。取 \(x\) 的基本邻域
\[
 W=\bigcap_{i\in F}\pi_i^{-1}(U_i),
\]
其中 \(F\) 有限且 \(U_i\) 是 \(x_i\) 的邻域。对每个 \(i\in F\)，存在阈值 \(\alpha_i\)，使
\(\alpha\ge\alpha_i\) 时 \(\pi_i(x_\alpha)\in U_i\)。有限次使用有向性得到共同后继 \(\alpha_0\)。
当 \(\alpha\ge\alpha_0\) 时所有有限坐标条件同时成立，故 \(x_\alpha\in W\)。每个邻域包含一个这样的
基本邻域，所以 \(x_\alpha\to x\)。
\end{proof}

\begin{proposition}[滤子的逐坐标收敛判别]\label{proposition:1-2-2-filter-coordinate}
乘积上的滤子 \(\mathcal F\) 收敛到 \(x=(x_i)\)，当且仅当
\((\pi_i)_*\mathcal F\to x_i\) 对每个 \(i\) 成立。
\end{proposition}
\begin{proof}
正向由坐标投影的连续性和滤子连续性判别得到。反过来，取 \(x\) 的基本柱邻域
\(W=\bigcap_{i\in F}\pi_i^{-1}(U_i)\)。坐标像滤子收敛给出
\(\pi_i^{-1}(U_i)\in\mathcal F\)。滤子对有限交封闭，故 \(W\in\mathcal F\)。因此每个 \(x\) 的邻域
属于 \(\mathcal F\)，即 \(\mathcal F\to x\)。
\end{proof}

\begin{example}[函数族的点态极限]\label{ex:1-2-2-1}
若每个 \(s\in S\) 的函数值都落在紧集 \(K_s\)，函数族 \(\mathcal A\) 可视为
\(\prod_sK_s\) 的子集。Tychonoff 定理将保证环境乘积紧，所以 \(\mathcal A\) 中任意网有逐点收敛子网。
极限先落在 \(\prod_sK_s\)；要使它仍属于 \(\mathcal A\)，还需 \(\mathcal A\) 对逐点极限封闭。
要把逐点收敛升级为一致收敛，则需要 \S\ref{subsec:1-3-2} 的等度连续性。

例如令 \(S=[0,1]\)、\(K_s=[0,1]\)，并令 \(\mathcal M\) 为所有单调不减映射
\(f:[0,1]\to[0,1]\) 的集合。若网 \(f_\alpha\in\mathcal M\) 逐点收敛到 \(f\)，则对
\(x\le y\) 有
\[
  f(x)=\lim_\alpha f_\alpha(x)\le \lim_\alpha f_\alpha(y)=f(y).
\]
所以 \(\mathcal M\) 在乘积中闭，从而在逐点拓扑下紧。这个结论并不声称极限连续；例如连续函数
\(f_n(x)=\min\{nx,1\}\) 逐点收敛到在 \(0\) 处不连续的单调函数。
\end{example}

\begin{example}[可数乘积度量的构造]\label{ex:1-2-2-2}
若 \(X_n\) 有度量 \(d_n\le1\)，则
\[
 d(x,y)=\sum_{n=1}^{\infty}2^{-n}d_n(x_n,y_n)
\]
定义一个度量：级数由 \(\sum_n2^{-n}=1\) 控制；若 \(d(x,y)=0\)，则每个非负项
\(2^{-n}d_n(x_n,y_n)\) 都为零，故 \(x_n=y_n\) 对所有 \(n\) 成立；三角不等式可对各坐标的
三角不等式逐项求和。\cref{proposition:1-2-2-3} 将进一步证明它诱导乘积拓扑。
\end{example}

\begin{example}[盒拓扑的双重失败]\label{ex:1-2-2-3}
若基本开集可以限制所有坐标，得到盒拓扑。在 \(\R^\N\) 中令
\(x^{(k)}_n=1/k\)。对每个固定 \(n\)，有 \(x^{(k)}_n\to0\)，所以它在乘积拓扑中趋于零；但它从不进入
盒邻域
\[
 \prod_{n=1}^{\infty}(-1/n,1/n),
\]
因为对任意 \(k\)，取 \(n>k\) 即有 \(1/k\notin(-1/n,1/n)\)。

盒拓扑也破坏紧性。考虑 \(D=\{(t,t,\ldots):t\in[0,1]\}\subset[0,1]^\N\)。若 \(x\notin D\)，
存在两个坐标 \(x_m\ne x_n\)。令 \(\eta=|x_m-x_n|/3\)，并在第 \(m,n\) 个坐标分别限制到
\[
  (x_m-\eta,x_m+\eta)\cap[0,1],\qquad
  (x_n-\eta,x_n+\eta)\cap[0,1].
\]
两区间不交；其余坐标取 \([0,1]\)，所得盒邻域不可能含对角点，故 \(D\) 闭。
对固定 \(t\)，在第 \(n\) 个坐标取相对邻域
\((t-1/n,t+1/n)\cap[0,1]\)。若对角点 \((s,s,\ldots)\) 属于这些区间的盒积，则
\[
  |s-t|<\frac1n\quad(n\ge1),
\]
从而 \(s=t\)。故该盒积与 \(D\) 只交于 \((t,t,\ldots)\)，所以 \(D\) 在子空间拓扑下离散。
单点族 \(\{\{p\}:p\in D\}\) 是 \(D\) 的开覆盖且没有有限子覆盖，故 \(D\) 不紧。若盒乘积紧，
其闭子集 \(D\) 应紧，矛盾。因此盒乘积不紧。
\end{example}

\begin{remark}[不可数乘积中的 Borel--柱差异]\label{rem:1-2-2-borel-cylinder}
为比较不可数乘积上的两种生成过程，在此使用以下集合概念：\(\sigma\)-代数是对补集与可数并封闭的
集合族，Borel \(\sigma\)-代数由全部开集生成。\S\ref{subsec:1-4-2} 将系统讨论这些定义。

令 \(I\) 不可数，\(X=\{0,1\}^{I}\)。称 \(E\subset X\) 依赖坐标集 \(J\subset I\)，若任意两点在
\(J\) 上坐标相同便同时属于或同时不属于 \(E\)。所有依赖某个可数坐标集的集合组成 \(\sigma\)-代数：
补集保留同一坐标集；可数并依赖对应可数坐标集之并。它包含全部有限柱，所以有限柱生成的乘积
\(\sigma\)-代数中，每个集合至多依赖可数多个坐标。

然而开集
\[
 X\setminus\{0^I\}=\bigcup_{i\in I}\{x:x_i=1\}
\]
不依赖任何可数 \(J\subset I\)：取 \(i\in I\setminus J\)，全零点与只在第 \(i\) 坐标为 \(1\) 的点在
\(J\) 上相同，却分别不属于和属于该开集。因此乘积拓扑的 Borel \(\sigma\)-代数可以严格大于有限柱生成的
\(\sigma\)-代数。这里使用“可数个可数集之并可数”时处在通常的 \textsf{CC}/ZFC 语境中。
\end{remark}

\begin{theorem}[Tychonoff 定理]\label{theorem:1-2-2-2}
若每个 \(X_i\) 紧，则 \(\prod_iX_i\) 在乘积拓扑下紧。
\end{theorem}
\begin{proof}
若某个因子为空，乘积为空，结论成立。以下设每个因子非空。

先在通常的 ZFC 语境中证明一般版本。给定乘积 \(X=\prod_iX_i\) 上的超滤子 \(\mathcal U\)，
\cref{proposition:1-1-3-4} 说明 \((\pi_i)_*\mathcal U\) 是 \(X_i\) 上超滤子。因 \(X_i\) 紧，
\cref{theorem:1-2-1-2} 给出非空极限集
\[
 L_i=\{x_i\in X_i:(\pi_i)_*\mathcal U\to x_i\}.
\]
对所有 \(i\) 同时选取 \(x_i\in L_i\)，得到 \(x=(x_i)\in X\)。由滤子的逐坐标判别，
\(\mathcal U\to x\)。所以 \(X\) 上每个超滤子都有极限，超滤子紧性定理给出 \(X\) 紧。这份证明既使用
超滤子扩张，也同时选择坐标极限，按 \textsf{AC} 记录。

若每个 \(X_i\) 还 Hausdorff，则每个 \(L_i\) 至多含一点；紧性又保证它非空，所以坐标极限由
\((\pi_i)_*\mathcal U\) 唯一确定，不需从一族多元素集合中同时选代表。在这种情形，采用
\textsf{UFL} 足以完成上述证明。本书后面使用的 \([0,1]^S\) 和紧值函数盒均属于紧 Hausdorff 情形。
\end{proof}

由 Tychonoff 定理，\cref{ex:1-2-1-1} 中的
\(\{0,1\}^{\mathcal P(\N)}\) 确实紧；结合该例中的坐标计算，它是一个紧而非序列紧的空间。

\begin{proposition}[可数乘积的第二可数性与可度量化]\label{proposition:1-2-2-3}
可数个第二可数空间的乘积第二可数；可数个可度量空间的乘积可度量。若每个因子紧且可度量，则乘积紧且
可度量，并可用对角子序列证明序列紧性。
\end{proposition}
\begin{proof}
对每个 \(n\)，取可数基 \(\mathcal B_n\)，并把 \(X_n\) 添入该基。所有只在某个有限坐标集
\(F\subset\N\) 上取 \(\mathcal B_n\) 成员、其余坐标不限制的柱集构成可数族；每个基本柱都含有
其中一个成员，所以乘积第二可数。

若 \(d_n\) 是 \(X_n\) 的度量，令 \(\rho_n=\min\{1,d_n\}\) 并定义
\[
 d(x,y)=\sum_{n=1}^{\infty}2^{-n}\rho_n(x_n,y_n).
\]
半径小于 \(1\) 的 \(\rho_n\)-球与同半径的 \(d_n\)-球相同，所以 \(\rho_n\) 与 \(d_n\) 诱导同一拓扑。
而且
\[
 \rho_n(x,z)\le \min\{1,d_n(x,y)+d_n(y,z)\}
 \le \rho_n(x,y)+\rho_n(y,z),
\]
所以 \(\rho_n\) 是度量，级数逐项给出 \(d\) 的三角不等式。若 \(d(x,y)=0\)，每项均为零，故 \(x=y\)。
固定 \(x\) 并给定 \(\varepsilon>0\)。选 \(N\) 使
\(\sum_{n>N}2^{-n}<\varepsilon/2\)，再令
\[
 W=\bigcap_{n=1}^N
   \pi_n^{-1}\!\left(B_{\rho_n}(x_n,\varepsilon/2)\right).
\]
若 \(y\in W\)，则前段和严格小于
\((\varepsilon/2)\sum_{n=1}^N2^{-n}<\varepsilon/2\)，尾段和也小于
\(\varepsilon/2\)，故 \(W\subset B_d(x,\varepsilon)\)。
反过来，设 \(x\) 的基本柱邻域只限制有限坐标集 \(F\subset\N\)。对每个 \(n\in F\)，取
\(r_n>0\) 使 \(B_{\rho_n}(x_n,r_n)\) 包含于相应坐标邻域。若 \(F\ne\varnothing\)，取
\[
 0<\delta<\min_{n\in F}2^{-n}\min\{1,r_n\}.
\]
若 \(d(x,y)<\delta\)，则对每个 \(n\in F\)，
\(2^{-n}\rho_n(x_n,y_n)<\delta\)，所以 \(y_n\) 落入相应坐标球。故 \(d\) 诱导乘积拓扑。
当 \(F=\varnothing\) 时柱邻域就是全空间，反向包含显然成立。

若每个因子紧可度量，Tychonoff 给出乘积紧。也可直接对任意序列逐坐标抽取：先取第一坐标收敛的子序列，
再从中取第二坐标收敛的子序列，如此递归；取第 \(k\) 层子序列的第 \(k\) 项作为对角序列。每个固定坐标
最终位于相应层子序列中，故逐坐标收敛，再由乘积度量收敛。递归抽取使用 \textsf{DC}/\textsf{CC}。
\end{proof}

\begin{proposition}[映入乘积的连续性]\label{proposition:1-2-2-4}
映射 \(F:Z\to\prod_iX_i\) 连续，当且仅当每个坐标映射 \(\pi_i\circ F\) 连续。
\end{proposition}
\begin{proof}
正向由连续映射复合得到。反过来，基本柱
\(W=\bigcap_{i\in E}\pi_i^{-1}(U_i)\) 的逆像为
\[
 F^{-1}(W)=\bigcap_{i\in E}(\pi_i\circ F)^{-1}(U_i),
\]
这是有限多个开集之交，因而开。基本柱的任意并给出全部开集，所以 \(F\) 连续。
\end{proof}

\begin{figure}[htbp]
\centering
\begin{tikzpicture}[node distance=1.35cm and 1.45cm,>=Stealth,every node/.style={align=center}]
 \node[book strong node] (u) {乘积超滤子 \(\mathcal U\)};
 \node[book node,right=of u] (p) {坐标超滤子\\\((\pi_i)_*\mathcal U\)};
 \node[book node,right=of p] (l) {坐标极限\\\(x_i\in L_i\)};
 \node[book node,below=of l] (x) {乘积点 \(x=(x_i)\)};
 \node[book node,left=of x] (c) {有限柱检验};
 \draw[book arrow] (u)--(p);
 \draw[book arrow] (p)--node[above]{因子紧}(l);
 \draw[book arrow] (l)--node[right]{一般版需选择}(x);
 \draw[book arrow] (u)|-(c);
 \draw[book arrow] (c)--node[below]{有限交}(x);
\end{tikzpicture}
\caption{坐标极限的存在与有限柱的收敛检验是两个不同步骤。}
\label{fig:1-2-2-tychonoff-coordinate-choice}
\end{figure}

乘积拓扑只给逐坐标控制；连续性和一致收敛需要额外结构。一般紧空间版、紧 Hausdorff 版和可数紧可度量版
的证明采用不同原则，不能把其中一份证明无条件搬到另一种因子类别。

\begin{exercise}
证明任意乘积中的坐标投影是开映射。
\end{exercise}
\begin{exercise}
在乘积非空的前提下，证明乘积 Hausdorff 当且仅当每个因子 Hausdorff；再用某个因子为空的情形说明为何
需要“乘积非空”。
\end{exercise}
\begin{exercise}
若每个 \((X_n,d_n)\) 完备，证明上述乘积度量完备。
\end{exercise}
\begin{exercise}
完整证明不可数 \(I\) 时，\(\{0,1\}^{I}\) 的有限柱生成 \(\sigma\)-代数严格小于乘积拓扑的 Borel
\(\sigma\)-代数。
\end{exercise}

乘积紧性负责产生候选极限，却不保证某个数值泛函在极限处保留正确的不等式。变分问题因此还需要一种
只控制“向下跳跃”的极限条件。

\subsection{半连续性、\texorpdfstring{\(\liminf\)}{liminf} 与紧性原理}
\label{subsec:1-2-3}

紧性能够阻止近似极小点逃走，却不能保证极限仍然极小。反过来，正确方向的极限不等式能够识别一个已经
存在的候选点，却不能自行产生候选。先看两种彼此独立的失败。

\begin{remark}
Weierstrass 极值定理通常为连续函数陈述，但直接法只需要禁止函数值在极限处向下跳，这正是下半连续。
把“候选点存在”与“候选点保持极小不等式”分开以后，紧性与半连续性的职责便不再混在一起。
\end{remark}

\begin{example}[两种最小化失败]\label{ex:1-2-3-3}
连续函数 \(f(x)=e^x\) 在 \(\R\) 上的下确界为 \(0\)，但没有最小点；点列趋向 \(-\infty\) 时函数值
趋零，候选点逃出了空间。第二种失败发生在紧集上：令
\[
 g(0)=1,\qquad g(x)=x\quad(0<x\le1).
\]
则 \(\inf_{[0,1]}g=0\)，仍没有最小点。这里空间不允许逃逸，失败来自
\(g(0)>\liminf_{x\downarrow0}g(x)\)。紧性与正确的极限不等式必须同时出现。
\end{example}

\begin{definition}[紧扩展实直线]\label{def:1-2-3-extended-line}
令 \(\overline\R=[-\infty,+\infty]\) 带次序拓扑。映射
\[
 h(t)=\frac{2}{\pi}\arctan t\quad(t\in\R),\qquad
 h(-\infty)=-1,\quad h(+\infty)=1
\]
是 \(\overline\R\) 到 \([-1,1]\) 的同胚。因此 \(\overline\R\) 是紧 Hausdorff 空间。
\end{definition}
\begin{proof}
\(h\) 在 \(\R\) 上严格递增、连续，并把 \(\R\) 双射到 \((-1,1)\)。当
\(t\to\pm\infty\) 时 \(h(t)\to\pm1\)，所以端点延拓连续。其逆在 \((-1,1)\) 上为
\(\tan(\pi s/2)\)，并在 \(s\to\pm1\) 时趋于相应无穷端点，故逆映射也连续。
\end{proof}

\begin{definition}[下半连续]\label{def:1-2-3-1}
函数 \(f:X\to\overline\R\) 称为\emph{下半连续}，若对每个 \(a\in\R\)，集合
\(\{x:f(x)>a\}\) 开。等价地，每个次水平集 \(\{x:f(x)\le a\}\) 闭。上半连续用相反不等式定义。
\end{definition}

\begin{proposition}[上图判别]\label{proposition:1-2-3-epigraph}
对 \(f:X\to\overline\R\)，令
\[
 \epi(f)=\{(x,t)\in X\times\R:f(x)\le t\}.
\]
则 \(f\) 下半连续，当且仅当 \(\epi(f)\) 在 \(X\times\R\) 中闭。
\end{proposition}
\begin{proof}
若上图闭，固定 \(a\in\R\)，连续映射 \(i_a:X\to X\times\R\)、\(i_a(x)=(x,a)\) 满足
\[
 i_a^{-1}(\epi(f))=\{x:f(x)\le a\},
\]
所以每个次水平集闭，\(f\) 下半连续。

反过来设 \(f\) 下半连续，取 \((x,t)\notin\epi(f)\)，即 \(f(x)>t\)。选实数 \(a\) 满足
\(t<a<f(x)\)。集合
\[
 \{y:f(y)>a\}\times(-\infty,a)
\]
是 \((x,t)\) 的开邻域；其中每个 \((y,s)\) 满足 \(f(y)>a>s\)，故不属于上图。因此上图的补集开，
上图闭。
\end{proof}

\begin{example}[边界值决定上图是否闭]\label{ex:1-2-3-step-epigraph}
在 \(\R\) 上，\(f=\1_{(0,\infty)}\) 下半连续：当 \(0\le a<1\) 时
\(\{f>a\}=(0,\infty)\)，其他阈值的超水平集是 \(\R\) 或空集。改令
\(g=\1_{[0,\infty)}\)，则 \(x_n=-1/n\to0\)，而
\[
 g(0)=1>\liminf_ng(x_n)=0.
\]
等价地，\((x_n,0)\in\epi(g)\)，其极限 \((0,0)\) 不在 \(\epi(g)\)。这个计算固定了下半连续不等式的
方向。
\end{example}

\begin{definition}[网的 \(\liminf\) 与 \(\limsup\)]\label{def:1-2-3-2}
对扩展实数网 \((a_\alpha)\)，定义
\[
 \liminf_\alpha a_\alpha
 =\sup_{\alpha_0}\inf_{\alpha\ge\alpha_0}a_\alpha,\qquad
 \limsup_\alpha a_\alpha
 =\inf_{\alpha_0}\sup_{\alpha\ge\alpha_0}a_\alpha.
\]
\end{definition}

\begin{proposition}[尾部聚点与 \(\liminf\)]\label{proposition:1-2-3-tail-cluster}
把 \((a_\alpha)\) 看成紧空间 \(\overline\R\) 中的网，并令
\[
 K_{\alpha_0}=\overline{\{a_\alpha:\alpha\ge\alpha_0\}},
 \qquad C=\bigcap_{\alpha_0}K_{\alpha_0}.
\]
则 \(C\) 是非空紧集，恰为该网的聚点集，而且
\[
 \min C=\liminf_\alpha a_\alpha,\qquad
 \max C=\limsup_\alpha a_\alpha.
\]
\end{proposition}
\begin{proof}
每个尾闭包都非空且闭。给定有限多个指标 \(\alpha_1,\ldots,\alpha_m\)，取共同后继
\(\gamma\)，则
\(K_\gamma\subset\bigcap_{j=1}^mK_{\alpha_j}\)；因此这些尾闭包具有有限交性质。扩展实直线紧，故
\(C\ne\varnothing\) 且紧。由
\cref{proposition:1-2-1-tail-closure}，\(C\) 正是聚点集。紧的非空扩展实数子集有最小值和最大值。
具体地，若非空闭集 \(K\subset\overline\R\) 且 \(b=\inf K\)，则 \(b=+\infty\) 时
\(K=\{+\infty\}\)；\(b=-\infty\) 时每个
\([-\infty,t)\) 都与 \(K\) 相交；\(b\in\R\) 时每个
\((b-\varepsilon,b+\varepsilon)\) 都与 \(K\) 相交。因此在所有情形都有 \(b\in\overline K=K\)；
上确界同理。

记 \(b_{\alpha_0}=\inf_{\alpha\ge\alpha_0}a_\alpha\) 和
\(\ell=\sup_{\alpha_0}b_{\alpha_0}\)。若
\(S_{\alpha_0}=\{a_\alpha:\alpha\ge\alpha_0\}\)，则
\(S_{\alpha_0}\subset[b_{\alpha_0},+\infty]\)。右侧闭，故
\(K_{\alpha_0}=\overline{S_{\alpha_0}}\subset[b_{\alpha_0},+\infty]\)，而
\(S_{\alpha_0}\subset K_{\alpha_0}\) 又给出反向的下确界不等式。因此闭包不改变下确界；上段论证还说明
该下确界属于闭包，故
\[
  \min K_{\alpha_0}=\min\overline{S_{\alpha_0}}
  =\inf S_{\alpha_0}=b_{\alpha_0}.
\]
所以 \(C\subset K_{\alpha_0}\) 给出 \(\min C\ge b_{\alpha_0}\)，进而
\(\min C\ge\ell\)。

若 \(\min C>\ell\)，取 \(t\) 使 \(\ell<t<\min C\)。由于
\(b_{\alpha_0}\le\ell<t\)，每个尾闭包都与闭区间 \([-\infty,t]\) 相交。更具体地，给定
有限多个指标 \(\alpha_1,\ldots,\alpha_m\)，取共同后继 \(\gamma\)；则
\(K_\gamma\subset\bigcap_jK_{\alpha_j}\)，而
\(\min K_\gamma=b_\gamma\le\ell<t\)，所以
\(K_\gamma\cap[-\infty,t]\ne\varnothing\)。这就证明闭集族
\[
 \{K_{\alpha_0}\}_{\alpha_0}\cup\{[-\infty,t]\}
\]
具有有限交性质；紧性给出 \(C\cap[-\infty,t]\ne\varnothing\)，与 \(\min C>t\) 矛盾。因此
\(\min C=\ell=\liminf_\alpha a_\alpha\)。最大值的论证把次序反向即可。
\end{proof}

\begin{definition}[次水平紧与等强制]\label{def:1-2-3-3}
函数 \(F:X\to\overline\R\) 称为\emph{次水平紧的}（inf-compact），若每个
\(\{F\le c\}\)、\(c\in\R\)，都紧。函数族 \((F_\lambda)\) 称为\emph{等强制的}，若对每个 \(c\in\R\)，
\[
 \bigcup_\lambda\{F_\lambda\le c\}
\]
相对紧。前者描述单个函数，后者统一控制一族函数。
\end{definition}

\begin{example}[可计算的等强制族]\label{ex:1-2-3-equicoercive}
在 \(\R\) 上令 \(F_n(x)=x^2+1/n\)。当 \(c<0\) 时所有次水平集均为空；当 \(c\ge0\) 时
\[
 \bigcup_n\{F_n\le c\}\subset[-\sqrt c,\sqrt c],
\]
所以 \((F_n)\) 等强制。若 \(\sup_nF_n(x_n)<\infty\)，则 \((x_n)\) 有界，因而相对紧。这个结论只产生
候选点；要比较变化中的 \(F_n\) 与某个极限泛函的值，还需额外的 \(\liminf\) 条件。相应的
\(\Gamma\)-收敛理论不在本节使用。
\end{example}

\begin{example}[开集指标与闭约束]\label{ex:1-2-3-1}
若 \(U\) 开，则 \(\1_U\) 下半连续；若 \(C\) 闭，则 \(\1_C\) 上半连续。对变分问题，更有用的闭约束
能量是
\[
 I_C(x)=
 \begin{cases}
  0,&x\in C,\\
  +\infty,&x\notin C.
 \end{cases}
\]
它的有限次水平集在 \(a\ge0\) 时为 \(C\)，在 \(a<0\) 时为空，所以 \(I_C\) 下半连续当且仅当 \(C\) 闭。
\end{example}

\begin{example}[有限维模型中的弱下半连续性]\label{ex:1-2-3-2}
先把本例所用术语就地定义：\(\R^d\) 的弱拓扑
\(\sigma(\R^d,(\R^d)^*)\) 是使每个线性泛函 \(\ell:\R^d\to\R\) 连续的最弱拓扑。记该拓扑为
\(\tau_w\)，欧氏拓扑为 \(\tau_e\)。每个线性泛函在 \(\tau_e\) 下连续，故
\(\tau_w\subset\tau_e\)。反过来，标准坐标泛函 \(p_j(x)=x_j\) 属于 \((\R^d)^*\)，而有限交
\[
 \bigcap_{j=1}^d p_j^{-1}\bigl((x_j-\varepsilon,x_j+\varepsilon)\bigr)
\]
构成欧氏拓扑的一套邻域基，故 \(\tau_e\subset\tau_w\)。于是 \(\tau_w=\tau_e\)。
Cauchy--Schwarz 不等式给出
\[
  \norm{x}=\sup_{\norm{u}\le1}\langle x,u\rangle,\qquad
  \{x:\norm{x}\le r\}
  =\bigcap_{\norm{u}\le1}\{x:\langle x,u\rangle\le r\},\qquad r\ge0.
\]
Cauchy--Schwarz 给出
\(\sup_{\norm{u}\le1}\langle x,u\rangle\le\norm{x}\)；当 \(x\ne0\) 时取
\(u=x/\norm{x}\) 即取到反向不等式，而 \(x=0\) 时等号显然。第二式于是由
\(\norm{x}\le r\) 当且仅当所有单位球内的 \(u\) 都满足
\(\langle x,u\rangle\le r\) 得到。
右端是闭半空间之交，所以闭球是弱闭集，范数因而弱下半连续。这个有限维论证只使用欧氏内积。
一般赋范空间中的相应闭半空间表示需要 Hahn--Banach 分离，见第~3~章。
\end{example}

\begin{theorem}[半连续性的网判别]\label{theorem:1-2-3-1}
函数 \(f:X\to\overline\R\) 下半连续，当且仅当对每个 \(x_\alpha\to x\)，
\[
 f(x)\le\liminf_\alpha f(x_\alpha).
\]
\end{theorem}
\begin{proof}
若 \(f(x)=-\infty\)，所需不等式没有内容。以下设 \(f(x)>-\infty\)。
先设 \(f\) 下半连续。给定 \(a<f(x)\)，开集 \(\{f>a\}\) 是 \(x\) 的邻域，所以存在
\(\alpha_0\) 使
\[
  \alpha\ge\alpha_0\quad\Longrightarrow\quad f(x_\alpha)>a.
\]
因此
\[
  \liminf_\alpha f(x_\alpha)
  =\sup_\beta\inf_{\alpha\ge\beta}f(x_\alpha)
  \ge \inf_{\alpha\ge\alpha_0}f(x_\alpha)\ge a.
\]
若 \(f(x)\in\R\)，令 \(a\uparrow f(x)\)；若 \(f(x)=+\infty\)，则上式对每个 \(a\in\R\)
成立。两种情形都给出所需不等式。

反过来，若 \(f\) 不下半连续，则存在 \(a\in\R\)，使 \(G=\{f>a\}\) 不是开集。取
\(x\in G\setminus\operatorname{int}G\)。每个 \(x\) 的邻域 \(U\) 都含有 \(y\) 满足 \(f(y)\le a\)。
令
\[
 J=\{(U,y):U\in\mathcal N(x),\ y\in U,\ f(y)\le a\},
\]
按邻域缩小排序，并令网值为 \(y\)。成对索引构造给出 \(y_j\to x\)。又因每个网值都满足
\(f(y_j)\le a\)，对每个阈值 \(j_0\) 都有
\[
  \inf_{j\ge j_0}f(y_j)\le a,
  \qquad
  \liminf_jf(y_j)=\sup_{j_0}\inf_{j\ge j_0}f(y_j)\le a<f(x),
\]
违背所假设的网不等式。
\end{proof}

\begin{remark}
若 \(X\) 第一可数，\cref{proposition:1-1-1-first-countable-sequential-tests} 中的邻域基构造把上一定理的
逆向反例网替换成序列，所以只检验序列即可。一般空间不能这样替换；本小节最后的习题给出反例。
\end{remark}

\begin{proposition}[上确界与有限和的稳定性]\label{proposition:1-2-3-2}
任意非空族 \(f_i:X\to\overline\R\) 的下半连续函数之逐点上确界仍下半连续。若
\(f_1,\ldots,f_n:X\to(-\infty,+\infty]\) 下半连续，则其逐点和也下半连续。
\end{proposition}
\begin{proof}
对上确界，
\[
 \left\{x:\sup_i f_i(x)>a\right\}
 =\bigcup_i\{x:f_i(x)>a\},
\]
右端开，故上确界下半连续。

只需证明两个函数之和。设 \(x_\alpha\to x\)。若 \(f(x),g(x)\) 都有限，给定 \(\varepsilon>0\)，
由下半连续性，
\[
 U_f=\{y:f(y)>f(x)-\varepsilon\},\qquad
 U_g=\{y:g(y)>g(x)-\varepsilon\}
\]
都是 \(x\) 的开邻域。收敛分别给出阈值 \(\alpha_f,\alpha_g\)，再由指标集的有向性取共同后继
\(\alpha_0\ge\alpha_f,\alpha_g\)。于是 \(\alpha\ge\alpha_0\) 时同时有
\[
 f(x_\alpha)>f(x)-\varepsilon,\qquad
 g(x_\alpha)>g(x)-\varepsilon.
\]
故 \(\liminf_\alpha(f+g)(x_\alpha)\ge f(x)+g(x)-2\varepsilon\)，再令
\(\varepsilon\downarrow0\)。若 \(f(x)=+\infty\)，则 \(g(x)>-\infty\)。给定 \(M\in\R\)，可取
实数 \(a,b\) 使 \(a+b=M\)、\(a<f(x)\) 且 \(b<g(x)\)：若 \(g(x)<+\infty\)，先取
\(b<g(x)\) 再令 \(a=M-b\)；若 \(g(x)=+\infty\)，可取 \(a=b=M/2\)。下半连续性给出共同尾部，
具体地，\(\{f>a\}\) 与 \(\{g>b\}\) 是 \(x\) 的开邻域；分别取收敛阈值并取二者的共同后继，
便使 \(f(x_\alpha)>a\) 且 \(g(x_\alpha)>b\)，从而和大于 \(M\)。因此和的 \(\liminf\) 为
\(+\infty\)。\(g(x)=+\infty\) 时同理。值域排除了 \(-\infty\)，所以没有
\(+\infty+(-\infty)\) 未定式。归纳得到有限和结论。
\end{proof}

\begin{theorem}[紧集上的极值]\label{theorem:1-2-3-3}
若 \(K\) 紧，\(f:K\to(-\infty,+\infty]\) 下半连续且不恒为 \(+\infty\)，则 \(f\) 在 \(K\) 上有限地
取到最小值。
\end{theorem}
\begin{proof}
令 \(m=\inf_Kf\)。先证 \(m>-\infty\)。若 \(m=-\infty\)，则闭集
\[
 F_n=\{x\in K:f(x)\le-n\}
\]
非空且递减，因而具有有限交性质。紧性给 \(x\in\bigcap_nF_n\)，这要求有限实数 \(f(x)\) 小于所有
\(-n\)，与值域矛盾。

因 \(f\) 不恒为 \(+\infty\)，\(m<+\infty\)。闭集
\[
 E_n=\{x\in K:f(x)\le m+1/n\}
\]
非空且递减。紧性给 \(x_*\in\bigcap_nE_n\)。于是 \(f(x_*)\le m+1/n\) 对所有 \(n\) 成立，故
\(f(x_*)\le m\)。下确界定义给相反不等式，因此 \(f(x_*)=m\in\R\)。
\end{proof}

\begin{proposition}[无选择的近似极小网]\label{proposition:1-2-3-minimizing-net}
设 \(F:X\to(-\infty,+\infty]\) 且 \(m=\inf_XF\in\R\)。存在一个网 \((x_j)\) 满足
\(F(x_j)\to m\)，其构造不需为每个误差预选一点。
\end{proposition}
\begin{proof}
令
\[
 J=\{(\varepsilon,x):\varepsilon>0,\ F(x)\le m+\varepsilon\},
\]
并规定
\[
 (\varepsilon,x)\preceq(\delta,y)\quad\Longleftrightarrow\quad\delta\le\varepsilon.
\]
下确界定义保证对每个 \(\varepsilon>0\) 都有可行点，所以 \(J\ne\varnothing\)。两个指标以
\(\eta=\min\{\varepsilon,\delta\}\) 加严，再取任一满足 \(F(z)\le m+\eta\) 的可行对，得到共同后继。
令网值 \(x_{(\varepsilon,x)}=x\)。给定 \(\rho>0\)，在任何误差坐标不超过 \(\rho\) 的指标之后，
\[
 m\le F(x_j)\le m+\rho.
\]
所以 \(F(x_j)\to m\)。指标集纳入全部可行对，没有指定整族 \(x_\varepsilon\)。
\end{proof}

\begin{theorem}[直接法原型]\label{theorem:1-2-3-4}
设 \(F:X\to(-\infty,+\infty]\) 下半连续，\(m=\inf_XF\in\R\)。若有网 \((x_\alpha)\) 满足
\(F(x_\alpha)\to m\)，且它的某个尾部落在一个相对紧集合中，则 \(F\) 取到 \(m\)。
\end{theorem}
\begin{proof}
设原网的指标集为 \(D\)。取阈值 \(\alpha_0\)，使尾网值落在相对紧集 \(A\) 中。闭包
\(K=\overline A\) 紧。把
\(D_0=\{\alpha\in D:\alpha\ge\alpha_0\}\) 作为尾指标集。它是有向集：\(D_0\) 中任意两个指标在原
指标集中有共同后继，而该后继仍大于 \(\alpha_0\)。包含映射 \(\iota:D_0\to D\) 最终共尾；事实上，
给定 \(\eta\in D\)，取 \(\eta\) 与 \(\alpha_0\) 的共同后继 \(\zeta\)，则
\(\alpha\ge\zeta\) 蕴含 \(\iota(\alpha)\ge\eta\)。因此尾网本身是原网的子网。

把该尾网看成 \(K\) 中的网；网紧性定理给出子网
\(x_{\iota(\phi(\beta))}\to x\in K\)，其中 \(\phi:B\to D_0\) 最终共尾。复合
\(\iota\circ\phi:B\to D\) 仍最终共尾：给定 \(\eta\in D\)，先取
\(\zeta\in D_0\) 使 \(\zeta\ge\eta\)，再在 \(\phi\) 的某个阈值后有
\(\phi(\beta)\ge\zeta\)。故这个收敛子网也是原网的子网，数值网
\(F(x_\alpha)\to m\) 沿它仍趋于 \(m\)。由下半连续性的网判别，
\[
 F(x)\le\liminf_\beta F(x_{\iota(\phi(\beta))})=m.
\]
另一方面 \(m\le F(x)\)，所以 \(F(x)=m\)。证明没有使用 Hausdorff 性；若空间非 Hausdorff，候选极限可能
不唯一，但任一由上述论证得到的极限仍是最小点。
\end{proof}

\begin{corollary}[次水平紧版直接法]\label{corollary:1-2-3-infcompact}
设 \(F:X\to(-\infty,+\infty]\) proper，即不恒为 \(+\infty\)，并且下半连续。若存在
\(c>\inf_XF\) 使 \(\{F\le c\}\) 相对紧，则 \(\inf_XF\in\R\) 且 \(F\) 取到它。特别地，proper、
下半连续且次水平紧的函数有最小点。
\end{corollary}
\begin{proof}
集合 \(E=\{F\le c\}\) 非空，并由下半连续性知其闭。相对紧性说明 \(\overline E\) 紧，而
\(E=\overline E\)，故 \(E\) 紧。紧集极值定理给出 \(x_*\in E\)，使 \(F(x_*)\) 是 \(F|_E\) 的有限
最小值。对 \(x\notin E\)，有 \(F(x)>c\ge F(x_*)\)，所以 \(x_*\) 也是全局最小点。这同时排除了
\(\inf_XF=-\infty\)。
\end{proof}

\begin{figure}[htbp]
\centering
\begin{tikzpicture}[node distance=1.2cm and 1.25cm,>=Stealth,every node/.style={align=center}]
 \node[book node] (m) {近似最小\\\(F(x_\alpha)\to m\)};
 \node[book node,below=of m] (s) {相对紧尾部};
 \node[book strong node,right=of m] (h) {同一网的\\两项信息};
 \node[book node,right=of h] (x) {收敛子网\\\(x_{\alpha_\beta}\to x\)};
 \node[book node,below=of x] (l) {下半连续\\\(F(x)\le m\)};
 \node[book strong node,left=of l] (e) {\(F(x)=m\)};
 \draw[book arrow] (m)--(h);
 \draw[book arrow] (s)--(h);
 \draw[book arrow] (h)--node[above]{紧性}(x);
 \draw[book arrow] (m)|-(l);
 \draw[book arrow] (x)--(l);
 \draw[book arrow] (l)--(e);
\end{tikzpicture}
\caption{紧性产生候选点，下半连续性识别候选点。}\label{fig:1-2-3-direct-method}
\end{figure}

下半连续本身既不保证下有界，也不保证取极值。任意族上半连续函数的上确界通常不再上半连续。等强制性只
控制候选点不逃逸，不自动给出变化泛函之间的极限不等式。第~3~章的 Fatou 引理和第~5~章的概率测度
开集不等式将在相应概念定义后给出同方向的 \(\liminf\) 机制；本章不使用这两个尚未陈述的结论。

\begin{exercise}
直接由定义证明：\(f\) 下半连续当且仅当每个次水平集闭。
\end{exercise}
\begin{exercise}
计算网 \(a_{(m,n)}=(-1)^m+1/n\)（\(\N^2\) 逐坐标排序）的 \(\liminf,\limsup\)，并求其聚点集。
\end{exercise}
\begin{exercise}
证明紧空间上的上半连续函数取到最大值。
\end{exercise}
\begin{exercise}
在通常的 ZFC 语境中，令 \(I=\mathcal P(\N)\)、\(X=\{0,1\}^{I}\)，并令
\[
 \Sigma=\{x\in X:\{i:x_i=1\}\text{ 至多可数}\}.
\]
证明 \(\Sigma\) 含所有有限支撑点，因而稠密但不闭；再证明 \(\Sigma\) 中任意逐坐标收敛序列的极限仍在
\(\Sigma\)。由此说明闭约束能量 \(I_\Sigma\) 满足所有序列的 \(\liminf\) 判别，却不下半连续。
\end{exercise}

下一节把紧性放入函数空间：点态候选只有在共同的局部控制下，才会成为连续并一致的函数极限。

\section{函数空间：从点态观测到一致控制}\label{sec:01-03}

拓扑紧性能够从许多有限条件中产生一个极限点，却不保证这个极限仍是连续函数。若
$f_\alpha(x)$ 对每个 $x$ 都收敛，极限函数甚至可能在每一点都与原函数族的整体形状毫无相似之处。
另一方面，只知道一族简单函数能区分任意两点，也还看不出它们为何能够逼近一个任意的连续函数。
本节处理这两道障碍。首先构造支撑可控的连续函数；随后把等度连续性解释为从有限采样向整个紧集传播
误差的机制；最后证明，函数代数的点分离能力经过紧性放大后会成为一致逼近能力。

\subsection{局部紧空间、紧化与连续函数工具}\label{subsec:1-3-1}

从一个很具体的任务开始。设 $K$ 是空间中的紧集，$U$ 是包含 $K$ 的开集。我们希望找到连续函数
$\varphi$，使它在 $K$ 上等于 $1$，在 $U$ 外等于 $0$，而且非零部分不会悄悄贴到 $U$ 的边界上。
在度量空间里，人们常写下距离函数来完成这件事；一般拓扑空间没有距离，必须找出距离公式真正用到的
分离性质。

\begin{example}[有限维局部紧性与 \(\ell^2\) 反例]\label{ex:1-3-1-1}
在 $\R^n$ 中，每个点 $x$ 都有闭包紧的球邻域。若 $G\subset\R^n$ 开且 $x\in G$，可取 $r>0$
使 $\overline{B(x,r)}\subset G$，因此 $G$ 也有同一性质；离散空间中，单点集就是所需的紧邻域。

这种现象在无限维赋范空间里已经会失败。以 $\ell^2$ 为例，若 $0$ 有一个闭包紧的邻域 $W$，则存在
$r>0$ 使 $B(0,r)\subset W$。闭球 $\overline{B(0,r/2)}$ 是 $\overline W$ 的闭子集，因而紧；
经线性缩放，$\ell^2$ 的闭单位球也应紧。然而标准正交向量满足
\[
  \norm{e_j-e_k}_2=\sqrt2\qquad(j\ne k),
\]
所以 $(e_k)$ 没有 Cauchy 子列，更没有收敛子列。这与紧度量空间的序列紧性矛盾。因此 $\ell^2$
在 $0$ 附近没有紧闭包的邻域。一般赋范空间何时局部紧是另一项结果；这里的计算已经提供了所需的
无限维反例。
\end{example}

\begin{definition}[局部紧 Hausdorff 空间]\label{def:1-3-1-1}
Hausdorff 空间 $X$ 称为\emph{局部紧的}，若每个 $x\in X$ 都有一个邻域，其在 $X$ 中的闭包紧。
记
\[
  A\Subset X
\]
表示 $\overline A$ 紧；写 $A\Subset U$ 时，还要求 $\overline A\subset U$。
\end{definition}

局部紧性允许空间整体向无穷远延伸，却要求每一点附近仍可用紧集控制。把所有逃向无穷远的方向合并成
一个点，会把许多局部构造化为紧空间上的分离问题。

\begin{definition}[一点紧化拓扑]\label{def:1-3-1-2}
设 $X$ 是非紧 Hausdorff 空间，在集合 $X^*=X\cup\{\infty\}$ 上规定：$X$ 中原有的开集仍然开；
含 $\infty$ 的集合 $W$ 为开集，当且仅当 $X^*\setminus W$ 是 $X$ 中的紧集。若这个拓扑是
Hausdorff 的，就称 $X^*$ 为 $X$ 的\emph{一点紧化}。
\end{definition}

定义中的条件确实给出一个拓扑。设一族规定的开集之并为 $O$，且其中某个成员 $W$ 含 $\infty$。令
$K=X^*\setminus W$，则 $K$ 是 $X$ 中的紧集。每个规定开集与 $X$ 的交都在 $X$ 中开：对含
$\infty$ 的成员，这是因为其补集紧且在 Hausdorff 空间 $X$ 中闭。因此 $O\cap X$ 在 $X$ 中开，且
\[
 X^*\setminus O=K\cap\bigl(X\setminus(O\cap X)\bigr)
\]
是 $K$ 的闭子集，因而紧。若没有成员含 $\infty$，该并就是 $X$ 中开集的并。有限交方面，
两个 $\infty$-邻域的交之补集是两个紧集之并；而 $X$ 中开集 $O$ 与
$\{\infty\}\cup(X\setminus K)$ 的交为 $O\setminus K$，它开是因为 Hausdorff 空间中的紧集闭。

\begin{proposition}[一点紧化判据]\label{proposition:1-3-1-one-point-compactification}
设 $X$ 是非紧 Hausdorff 空间，并在 $X^*$ 上采用\cref{def:1-3-1-2}中的拓扑。则 $X^*$ 紧，
$X$ 是 $X^*$ 的开稠密子空间，并且
\[
  X^*\text{ 是 Hausdorff 空间}
  \quad\Longleftrightarrow\quad
  X\text{ 是局部紧空间}.
\]
\end{proposition}

\begin{proof}
给定 $X^*$ 的开覆盖，选取其中一个含 $\infty$ 的成员 $W$。集合
$K=X^*\setminus W$ 是 $X$ 中的紧集，故覆盖中有限多个成员已经覆盖 $K$；再添 $W$ 就得到
$X^*$ 的有限子覆盖。于是 $X^*$ 紧。$X$ 是开子空间，并保持原拓扑。又因 $X$ 非紧，每个形如
$\{\infty\}\cup(X\setminus K)$ 的邻域都与 $X$ 相交，所以 $\infty\in\overline X$，即 $X$ 稠密。

先设 $X$ 局部紧。$X$ 内的两点可由原来的 Hausdorff 性分离。给定 $x\in X$，取开邻域 $V$
使 $\overline V$ 紧。于是
\[
  V
  \quad\text{与}\quad
  \{\infty\}\cup(X\setminus\overline V)
\]
是分别包含 $x$ 与 $\infty$ 的不交开集，故 $X^*$ 为 Hausdorff 空间。

反过来设 $X^*$ 为 Hausdorff 空间。给定 $x\in X$，可取不交开集 $V,W\subset X^*$，分别包含
$x$ 与 $\infty$。按拓扑的定义，存在紧集 $K\subset X$ 使
$\{\infty\}\cup(X\setminus K)\subset W$。因此 $V\cap X\subset K$。紧集 $K$ 在 Hausdorff
空间 $X$ 中闭，故 $\overline{V\cap X}^{\,X}\subset K$，这个闭包是紧集。于是 $x$ 有闭包紧的
邻域 $V\cap X$，从而 $X$ 局部紧。
\end{proof}

\begin{definition}[支撑与三类连续函数]\label{def:1-3-1-3}
标量域 $\mathbb F$ 固定为 $\R$ 或 $\C$。对 $f\in C(X,\mathbb F)$，定义
\[
  \supp f=\overline{\{x\in X:f(x)\ne0\}}.
\]
记 $C_b(X)$ 为所有有界连续标量函数组成的空间，$C_c(X)$ 为其中支撑紧的函数组成的子空间。
函数 $f\in C(X,\mathbb F)$ 称为\emph{在无穷远消失}，若对每个 $\varepsilon>0$，集合
\[
  \{x\in X:|f(x)|\geq\varepsilon\}
\]
紧；所有这类函数组成 $C_0(X)$。三个空间均使用上确界范数
$\norm{f}_\infty=\sup_{x\in X}|f(x)|$；当 $X=\varnothing$ 时，约定唯一函数的上确界范数为 $0$。
上述条件蕴含 $f$ 有界，因而这个范数在 $C_0(X)$ 上有限。
\end{definition}

最后一句值得核对。取 $\varepsilon=1$，集合 $K=\{|f|\ge1\}$ 紧，故 $f|_K$ 有界；在 $X\setminus K$
上已有 $|f|<1$。于是 $f$ 在整个 $X$ 上有界。若 $X$ 非紧且局部紧，则
\[
 f\in C_0(X)
 \quad\Longleftrightarrow\quad
 \widetilde f(x)=
 \begin{cases}
   f(x),&x\in X,\\
   0,&x=\infty
 \end{cases}
 \text{ 属于 }C(X^*).
\]
事实上，$\widetilde f$ 在 $\infty$ 连续恰好表示：对每个 $\varepsilon>0$，存在紧集 $K$，使
$|f|<\varepsilon$ 于 $X\setminus K$。若 $f\in C_0(X)$，可直接取
$K_\varepsilon=\{|f|\ge\varepsilon\}$。反过来，若存在这样的紧集 $K$，则
\[
  \{|f|\ge\varepsilon\}\subset K.
\]
左侧由 $f$ 的连续性在 $X$ 中闭，因而是紧集 $K$ 的闭子集，所以 $f\in C_0(X)$。若 $X$ 紧，则每个
连续函数都在无穷远消失，故
$C_0(X)=C(X)$。

\begin{example}[无穷远处的连续性]\label{ex:1-3-1-2}
在 $\R^n$ 上令
\[
  f(x)=\frac{1}{1+\norm{x}}.
\]
当 $0<\varepsilon\le1$ 时，
$\{f\ge\varepsilon\}=\{\norm{x}\le\varepsilon^{-1}-1\}$ 是闭有界集；当 $\varepsilon>1$ 时该集合为空。
所以 $f\in C_0(\R^n)$，其在一点紧化上满足 $f(\infty)=0$。常数函数 $1$ 则不在
$C_0(\R^n)$ 中，因为 $\{|1|\ge1/2\}=\R^n$ 不紧。这个差别不是关于函数是否有界，而是关于它在
逃离每个紧集后是否一致趋于零。
\end{example}

接下来先在紧 Hausdorff 空间上建立所需的分离工具。证明包含两个层次：拓扑的正规性给出嵌套开集，
二进有理数把这些开集编码成一个连续实函数。

\begin{theorem}[紧 Hausdorff 连续分离工具包]\label{theorem:1-3-1-compact-separation}
每个紧 Hausdorff 空间 $L$ 都是正规的，即任意两个不交闭集都可由不交开集分离，并且有以下结论。
\begin{enumerate}
  \item 若闭集 $A,B\subset L$ 不交，则存在 $u\in C(L,[0,1])$，使
        $u|_A=0$ 且 $u|_B=1$。
  \item 若 $A\subset L$ 闭，$M\ge0$，且 $f\in C(A,[-M,M])$，则存在
        $F\in C(L,[-M,M])$ 使 $F|_A=f$。
\end{enumerate}
\end{theorem}

\begin{proof}
先证正规性。给定 $x\notin B$，考虑所有满足下列条件的开集 $Q$：存在包含 $x$ 的开集 $P$，
使 $P\cap Q=\varnothing$。Hausdorff 性说明这些 $Q$ 覆盖 $B$。闭集 $B$ 是紧集，故其中有限多个
$Q_1,\ldots,Q_m$ 覆盖 $B$；对这有限多个 $Q_j$ 分别取相应的 $P_j$。于是
\[
  P=\bigcap_{j=1}^mP_j,\qquad Q=\bigcup_{j=1}^mQ_j
\]
是不交开集，分别包含 $x$ 与 $B$。现在令 $A,B$ 为不交闭集。对每个 $x\in A$ 作上述点与闭集的
分离。所有能够与某个包含 $B$ 的开集分离的开集 $P$ 构成 $A$ 的开覆盖。由 $A$ 的紧性取其中有限个
$P_1,\ldots,P_r$ 覆盖 $A$，并对这有限多个 $P_i$ 取相应的 $Q_i\supset B$；则
\[
  \bigcup_{i=1}^rP_i
  \quad\text{与}\quad
  \bigcap_{i=1}^rQ_i
\]
给出所需分离。

正规性还有一个反复使用的等价形式：若闭集 $F\subset O$ 且 $O$ 开，则存在开集 $V$ 使
\begin{equation}\label{eq:1-3-normal-shrink}
  F\subset V\subset\overline V\subset O.
\end{equation}
为此分离 $F$ 与闭集 $L\setminus O$，设相应不交开集为 $V$ 与 $W$。由
$\overline V\subset L\setminus W\subset O$ 即得\eqref{eq:1-3-normal-shrink}。

现在证明第一项。令 $D=\{k/2^n:0\le k\le2^n,\ n\in\N\}$。先置
$U_1=L\setminus B$，再由\eqref{eq:1-3-normal-shrink}取 $U_0$，使
$A\subset U_0\subset\overline U_0\subset U_1$。逐层插入奇数分点：若第 $n$ 层的开集已经选定，
就在每一对相邻的 $U_{k/2^n}$ 与 $U_{(k+1)/2^n}$ 之间使用
\eqref{eq:1-3-normal-shrink}，选取 $U_{(2k+1)/2^{n+1}}$，使
\[
 \overline U_{k/2^n}
 \subset U_{(2k+1)/2^{n+1}}
 \subset\overline U_{(2k+1)/2^{n+1}}
 \subset U_{(k+1)/2^n}.
\]
于是对任意 $r,s\in D$ 且 $r<s$，都有 $\overline U_r\subset U_s$。定义
\[
  u(x)=\inf\{r\in D:x\in U_r\},
\]
并在右侧集合为空时令 $u(x)=1$。$A\subset U_0$ 给出 $u|_A=0$，而 $B$ 与每个 $U_r$
都不交，故 $u|_B=1$。

只余连续性。对 $0<a<1$，稠密性与定义给出
\[
  \{u<a\}=\bigcup_{\substack{r\in D\\r<a}}U_r.
\]
另一方面，
\[
  \{u>a\}=\bigcup_{\substack{r\in D\\r>a}}(L\setminus\overline U_r).
\]
为验证第二式的正向包含，若 $u(x)>a$，选 $a<r<s<u(x)$；此时 $x\notin U_s$，而
$\overline U_r\subset U_s$，所以 $x\notin\overline U_r$。反向包含则由
$x\notin\overline U_r$ 推出 $x\notin U_t$ 对每个 $t<r$，因而 $u(x)\ge r>a$。
两个集合都是开集；对阈值 $a\le0$ 或 $a\ge1$，严格次水平集与超水平集由 $0\le u\le1$
直接处理。所以 $u$ 连续。这个二进递归只进行可数多层选择；用后文的记号，它采用
$\textsf{DC}$。

证明第二项。$M=0$ 时取零函数。设 $M>0$，先以 $M$ 缩放，归结为 $|f|\le1$。我们使用如下
一次逼近。若 $r\in C(A)$ 且 $|r|\le q$，令
\[
 A_- =\{r\le-q/3\},\qquad A_+=\{r\ge q/3\}.
\]
这两个闭集不交。由第一项取 $v\in C(L,[0,1])$，使 $v=0$ 于 $A_-$、$v=1$ 于 $A_+$，并令
$g=(q/3)(2v-1)$。于是 $\norm{g}_\infty\le q/3$。在 $A_+$ 上，$g=q/3$；在 $A_-$ 上，
$g=-q/3$；在余下点，$|r|<q/3$ 且 $|g|\le q/3$。逐一比较得到
\begin{equation}\label{eq:1-3-tietze-step}
  |r-g|\le 2q/3\qquad\text{于 }A.
\end{equation}

令 $r_0=f$。假设 $r_n\in C(A)$ 已经定义且
$\norm{r_n}_{C(A)}\le(2/3)^n$。把一次逼近用于 $q=(2/3)^n$，选取
$g_n\in C(L)$，再定义
\[
  r_{n+1}=r_n-g_n|_A
  =f-\sum_{k=0}^{n}(g_k|_A).
\]
由\eqref{eq:1-3-tietze-step}，
$\norm{r_{n+1}}_{C(A)}\le(2/3)^{n+1}$，所以归纳确实可以继续，并得到
\[
 \norm{g_n}_\infty\le\frac13\left(\frac23\right)^n,
 \qquad
 \left|f-\sum_{k=0}^{n}g_k\right|
 \le\left(\frac23\right)^{n+1}\quad\text{于 }A.
\]
级数 $\sum_ng_n$ 在 $L$ 上一致收敛为连续函数 $F$；第一个估计给出
$\norm{F}_\infty\le\sum_{n\ge0}\frac13(2/3)^n=1$，第二个估计给出 $F|_A=f$。
最后乘回 $M$，即得所述延拓。这里逐次选取残差逼近函数也采用可数递归选择，集合论标记同为
$\textsf{DC}$。
\end{proof}

有限离散空间为这个构造提供了直接检验：Urysohn 函数就是闭集的 $0/1$ 分配，Tietze 延拓只是在未指定
点选取区间内的数值。二进开集族的作用，是把同一分配改造成连续函数。需要强调的是，一般 Hausdorff
空间未必正规。一个标准反例是 Sorgenfrey 直线 $S$ 的平方，其中 $S$ 的拓扑基由半开区间
$[a,b)$ 组成。若 $x<y$，则 $[x,y)$ 与 $[y,y+1)$ 是不交基本邻域，故 $S$ 为 Hausdorff 空间，
从而 $S^2$ 也 Hausdorff。每个非空基本区间 $[a,b)$ 都含有有理数
$q$（取 $a\le q<b$），所以 $\Q$ 在 $S$ 中稠密，进而 $\Q^2$ 是 $S^2$ 的可数稠密子集。令
\[
 D=\{(x,-x):x\in\R\}
\]
是基数为连续统的闭离散子空间。若 $(a,b)\notin D$，则 $a+b\ne0$：当 $a+b>0$ 时，任意
$\varepsilon>0$ 给出的矩形 $[a,a+\varepsilon)\times[b,b+\varepsilon)$ 内各点坐标和都为正；当
$a+b<0$ 时取 $0<\varepsilon<-(a+b)/2$，矩形内各点坐标和都小于零。故 $D$ 闭。又若
$(y,-y)\in[x,x+\varepsilon)\times[-x,-x+\varepsilon)$，则同时有 $y\ge x$ 与 $y\le x$，所以
$y=x$；故该矩形与 $D$ 只交于 $(x,-x)$，从而 $D$ 离散。若 $S^2$ 正规，则对
$D$ 的每个子集 $E$，离散性说明 $E$ 与 $D\setminus E$ 都在 $D$ 中闭；再由 $D$ 在 $S^2$ 中闭，
它们也在 $S^2$ 中闭。上面 Urysohn 证明的二进开集构造在正规性建立后只使用正规性，不再使用紧性，
故在这个反设下可照搬，并给出一个在 $E$ 上取 $0$、在 $D\setminus E$ 上取 $1$ 的
连续函数。下面把所需的基数矛盾写成显式映射。记
$\mathscr C=C(S^2,\R)$，并定义
\[
 \Theta:\mathscr C\longrightarrow\mathcal P(D),\qquad
 \Theta(f)=\{d\in D:f(d)<1/2\}.
\]
若 $S^2$ 正规，则对每个 $E\subset D$，上述分离函数 $f$ 满足 $\Theta(f)=E$，所以 $\Theta$ 满射。

另一方面，限制映射
\[
 R:\mathscr C\longrightarrow\R^{\Q^2},\qquad R(f)=f|_{\Q^2}
\]
是单射。事实上，若 $f|_{\Q^2}=g|_{\Q^2}$ 而某点 $x$ 满足 $f(x)\ne g(x)$，则连续函数
$|f-g|$ 的开集
\[
 O=\{z\in S^2:|f(z)-g(z)|>\tfrac12|f(x)-g(x)|\}
\]
是 $x$ 的邻域；它与稠密集 $\Q^2$ 相交，和限制相等矛盾。
还可把 $\R^{\Q^2}$ 单射编码进 $\R$：固定枚举
$\Q^2=\{z_n:n\ge1\}$、$\Q=\{q_m:m\ge1\}$ 以及双射
$\pi:\N^2\to\N$。对数列 $a=(a_n)$ 令
\[
 E_a=\{(n,m):q_m<a_n\},\qquad
 c(a)=\sum_{(n,m)\in E_a}2\,3^{-\pi(n,m)}.
\]
有理数的稠密性说明 $a\ne b$ 时 $E_a\ne E_b$；而三进制数位只取 $0,2$，故上式把不同的
$E_a$ 送到不同实数。更直接地，若 $k$ 是两个编码数位的最小差异位置，则该位贡献
$2\cdot3^{-k}$，其后全部数位贡献的绝对值至多
$\sum_{j>k}2\cdot3^{-j}=3^{-k}$，所以差不为零。因此 $c$ 是单射，进而 $\mathscr C$ 可单射进
$\R$。又 $x\mapsto(x,-x)$ 是 $\R$ 到 $D$ 的双射，故存在单射 $j:\mathscr C\to D$。

常值零函数 $\mathbf0$ 属于 $\mathscr C$。利用它定义满射
\[
 p:D\longrightarrow\mathscr C,\qquad
 p(d)=
 \begin{cases}
   j^{-1}(d),&d\in j(\mathscr C),\\
   \mathbf0,&d\notin j(\mathscr C).
 \end{cases}
\]
于是 $\Theta\circ p:D\to\mathcal P(D)$ 满射。但对任何映射
$h:D\to\mathcal P(D)$，对角集
$H=\{d\in D:d\notin h(d)\}$ 都不可能等于任一 $h(d)$，所以这样的满射不存在。这一矛盾证明
$S^2$ 不正规，也说明不能在任意 Hausdorff 空间中无条件使用连续分离工具。

\begin{proposition}[相对紧邻域引理]\label{proposition:1-3-1-1}
设 $X$ 局部紧 Hausdorff，$K\subset U\subset X$，其中 $K$ 紧、$U$ 开。则存在开集 $V$ 使
\[
  K\subset V\Subset U.
\]
\end{proposition}

\begin{proof}
若 $X$ 非紧，就在紧 Hausdorff 空间 $X^*$ 中工作；若 $X$ 紧，就在 $X$ 中工作。记这个紧空间为
$L$。所有满足 $\overline W^{\,L}\subset U$ 的 $L$-开集 $W$ 覆盖 $K$：对每个 $x\in K$，
把\eqref{eq:1-3-normal-shrink}用于闭集 $\{x\}\subset U$ 即可。紧性给出其中有限个
$W_1,\ldots,W_m$ 覆盖 $K$。令
$V=X\cap\bigcup_{j=1}^mW_j$。它是 $X$ 中的开集，包含 $K$，而
\[
 \overline V^{\,X}\subset\bigcup_{j=1}^m\overline{W_j}^{\,L}\subset U.
\]
右侧有限并是 $L$ 中的紧集且包含于 $X$，所以也是 $X$ 中的紧集。故 $V\Subset U$。
\end{proof}

\begin{theorem}[局部 Urysohn 引理]\label{theorem:1-3-1-2}
设 $X$ 局部紧 Hausdorff，$K\subset U\subset X$，其中 $K$ 紧、$U$ 开。则存在
$\varphi\in C_c(X,[0,1])$，使
\[
  \varphi|_K=1,\qquad \supp\varphi\subset U.
\]
\end{theorem}

\begin{proof}
由\cref{proposition:1-3-1-1}取开集 $V$ 使 $K\subset V\Subset U$。若 $X$ 非紧，在 $X^*$
中对不交闭集 $X^*\setminus V$ 与 $K$（按此顺序）使用
\cref{theorem:1-3-1-compact-separation}；若 $X$ 紧，就在 $X$ 中对 $X\setminus V$ 与 $K$
按同一顺序分离。由该定理中“第一闭集取值 $0$、第二闭集取值 $1$”的约定，所得连续函数
$\varphi$ 在 $K$ 上为 $1$，在 $X\setminus V$ 上为 $0$。
限制到 $X$ 后仍记作 $\varphi$。非零集包含于 $V$，故
$\supp\varphi\subset\overline V\subset U$；又 $\overline V$ 紧，所以 $\varphi\in C_c(X)$。
\end{proof}

\begin{example}[支撑受控的截断]\label{ex:1-3-1-3}
在 $\R$ 上取 $K=[-1,1]$、$U=(-2,2)$，局部 Urysohn 引理给出所需截断。这里还能写出一个公式：
\[
 \varphi(x)=
 \begin{cases}
  1,&|x|\le1,\\
  2-|x|,&1<|x|<2,\\
  0,&|x|\ge2.
 \end{cases}
\]
它连续，$\varphi|_K=1$，且 $\supp\varphi=[-2,2]\subset U$ 这一句其实不成立：端点
$\pm2$ 不属于 $U$。这正暴露了手写公式常见的边界疏忽。把中间断点改为 $3/2$，即令第二行在
$1<|x|<3/2$ 上取 $3-2|x|$、在 $|x|\ge3/2$ 时取零，才得到
$\supp\varphi=[-3/2,3/2]\subset U$。抽象定理把这种严格的支撑余量直接写进结论。
\end{example}

\begin{figure}[htbp]
\centering
\begin{tikzpicture}[scale=0.92]
  \draw[book region,fill=BookAccent!3] (0,0) ellipse (5.2cm and 2.4cm);
  \draw[book region,fill=BookAccent!6] (0,0) ellipse (3.9cm and 1.75cm);
  \draw[book region,fill=BookAccent!10] (0,0) ellipse (2.5cm and 1.05cm);
  \draw[book region,fill=BookWarm!12,draw=BookWarm] (0,0) ellipse (1.15cm and 0.45cm);
  \node[book label] at (0,0) {$K$，$\varphi=1$};
  \node[book label] at (0,0.76) {$V$};
  \node[book label] at (0,1.47) {$\supp\varphi\subset U$};
  \node[book label] at (0,2.14) {$X$};
  \node[book label,anchor=west] at (4.1,1.25) {$\varphi=0$};
  \draw[book accent arrow] (1.2,-0.2) -- (2.4,-0.65)
    node[midway,below,book label] {$K\subset V\Subset U$};
\end{tikzpicture}
\caption{相对紧邻域给截断函数留下严格的支撑余量；局部 Urysohn 引理在这些嵌套集合之间构造截断函数。}
\label{fig:1-3-1-cutoff}
\end{figure}

\begin{theorem}[Tietze 延拓的紧支撑版本]\label{theorem:1-3-1-3}
设 $X$ 局部紧 Hausdorff，$K\subset U\subset X$，其中 $K$ 紧、$U$ 开。若
$f\in C(K,\mathbb F)$，则存在 $\widetilde f\in C_c(X,\mathbb F)$，满足
\[
 \widetilde f|_K=f,\qquad
 \supp\widetilde f\subset U,\qquad
 \norm{\widetilde f}_\infty=\norm f_\infty.
\]
\end{theorem}

\begin{proof}
令 $M=\norm f_\infty$。若 $M=0$，取 $\widetilde f=0$。以下设 $M>0$。由
\cref{proposition:1-3-1-1}取 $K\subset V\Subset U$。在紧 Hausdorff 空间 $\overline V$ 上，
实值情形由\cref{theorem:1-3-1-compact-separation}得到延拓 $F$，满足
$F|_K=f$ 且 $|F|\le M$。

复值情形需要多一步。分别延拓 $\Re f$ 与 $\Im f$，得到连续映射
$G:\overline V\to\R^2\cong\C$。定义闭圆盘上的径向收缩
\[
 R_M(z)=
 \begin{cases}
 z,&|z|\le M,\\
 Mz/|z|,&|z|>M.
 \end{cases}
\]
$R_M$ 连续，且在半径 $M$ 的闭圆盘上为恒等映射。因此 $F=R_M\circ G$ 仍延拓 $f$，并满足
$|F|\le M$。这一步说明，分别延拓实部与虚部后不能直接宣称复模长不增；径向收缩才恢复精确范数界。

再由\cref{theorem:1-3-1-2}取 $\psi\in C_c(X,[0,1])$，使 $\psi=1$ 于 $K$ 且
$\supp\psi\subset V$。在 $X$ 上定义
\[
 \widetilde f(x)=
 \begin{cases}
   \psi(x)F(x),&x\in\overline V,\\
   0,&x\notin\overline V.
 \end{cases}
\]
因为 $\supp\psi\subset V$，若 $x\in\partial V$，则
$x\notin\supp\psi$；闭集 $\supp\psi$ 的补集给出 $x$ 的开邻域 $N_x$，且
\[
  \psi|_{N_x}=0.
\]
故分段函数在 $N_x$ 上恒为零，在每个边界点连续；在 $\overline V$ 的内部及其补集内部，连续性直接来自
各自的定义。
它在 $K$ 上等于 $f$，支撑包含于 $\supp\psi\subset V\subset U$，并且
$\norm{\widetilde f}_\infty\le M$。另一方面，限制到 $K$ 给出
$\norm{\widetilde f}_\infty\ge\norm f_\infty=M$，故等号成立。
\end{proof}

\begin{proposition}[有限单位分解]\label{proposition:1-3-1-finite-partition-of-unity}
设 $X$ 局部紧 Hausdorff，紧集 $K$ 被开集 $U_1,\ldots,U_m$ 覆盖。则存在
$\varphi_j\in C_c(X,[0,1])$，使
\[
  \supp\varphi_j\subset U_j\quad(1\le j\le m),
  \qquad
  \sum_{j=1}^m\varphi_j=1\quad\text{于 }K.
\]
\end{proposition}

\begin{proof}
考虑所有满足
\[
 W\text{ 开},\qquad \overline W\text{ 紧},\qquad
 \overline W\subset U_j\text{ 对某个 }j
\]
的开集 $W$。由\cref{proposition:1-3-1-1}，它们覆盖 $K$。紧性给出有限子覆盖
$W_1,\ldots,W_N$。对每个 $i$ 选一个 $j_i$ 使 $\overline W_i\subset U_{j_i}$；这里只作有限次选择。
由局部 Urysohn 引理，存在 $\psi_i\in C_c(X,[0,1])$，使
$\psi_i=1$ 于 $\overline W_i$ 且 $\supp\psi_i\subset U_{j_i}$。令
\[
  \psi_j=\sum_{\{i:j_i=j\}}\psi_i,
  \qquad s=\sum_{j=1}^m\psi_j.
\]
空和取零。则 $s>0$ 于 $K$，且每个 $\psi_j$ 有紧支撑并满足
$\supp\psi_j\subset U_j$。令 $O=\{s>0\}$。再用局部 Urysohn 引理取
$\chi\in C_c(X,[0,1])$，使 $\chi=1$ 于 $K$ 且 $\supp\chi\subset O$。定义
\[
 \varphi_j(x)=
 \begin{cases}
   \chi(x)\psi_j(x)/s(x),&x\in O,\\
   0,&x\notin O.
 \end{cases}
\]
由于 $\supp\chi$ 是 $O$ 的闭子集，分段定义在 $\partial O$ 附近恒为零，所以 $\varphi_j$ 连续。
在 $O$ 上有 $0\le\psi_j\le s$，所以 $0\le\varphi_j\le\chi\le1$；并有
$\sum_j\varphi_j=\chi=1$ 于 $K$。又
$\supp\varphi_j\subset\supp\psi_j\subset U_j$，结论成立。
\end{proof}

\begin{proposition}[$C_c$ 在 $C_0$ 中稠密]\label{proposition:1-3-1-4}
设 $X$ 局部紧 Hausdorff。则 $C_c(X)$ 在 $C_0(X)$ 中关于上确界范数稠密。
\end{proposition}

\begin{proof}
给定 $f\in C_0(X)$ 与 $\varepsilon>0$，集合
$K=\{x:|f(x)|\ge\varepsilon/2\}$ 紧。对 $K\subset X$ 用局部 Urysohn 引理取
$\varphi\in C_c(X,[0,1])$，使 $\varphi=1$ 于 $K$。于是 $\varphi f\in C_c(X)$。在 $K$ 上
$f-\varphi f=0$；在 $X\setminus K$ 上，
\[
 |f-\varphi f|=|1-\varphi|\,|f|<\varepsilon/2.
\]
所以 $\norm{f-\varphi f}_\infty\le\varepsilon/2<\varepsilon$，这就给出稠密性。
\end{proof}

对任意开集 $U\subset X$，局部 Urysohn 引理还给出一个精确的点态公式：
\[
 \mathbf1_U(x)=
 \sup\{\varphi(x):\varphi\in C_c(X,[0,1]),\ \supp\varphi\subset U\}.
\]
若 $x\notin U$，右侧每一项都为零；若 $x\in U$，对紧集 $\{x\}\subset U$ 应用该引理，右侧就达到
$1$。以后讨论 Radon 正则性时，这个公式会把开集的示性函数转化为支撑受控的连续测试函数。

\begin{proposition}[$C_0(X)$ 的完备性]\label{proposition:1-3-1-c0-complete}
设 $X$ 为拓扑空间，标量域为 $\R$ 或 $\C$。则 $C_0(X)$ 在上确界范数下是 Banach 空间。
\end{proposition}

\begin{proof}
设 $(f_n)$ 是 $C_0(X)$ 中的 Cauchy 列。对每个 $x\in X$，标量列 $(f_n(x))$ 是 Cauchy 列，记其
极限为 $f(x)$。给定 $\varepsilon>0$，取 $N$ 使
$\norm{f_n-f_m}_\infty<\varepsilon$ 对所有 $m,n\ge N$ 成立；固定 $n\ge N$ 并令
$m\to\infty$，得到 $|f_n(x)-f(x)|\le\varepsilon$ 对每个 $x$ 成立。因此 $f_n\to f$ 一致，
$f$ 连续且有界。

还须核对无穷远消失。固定 $\varepsilon>0$，取 $n$ 使
$\norm{f-f_n}_\infty<\varepsilon/2$。于是
\[
 \{x:|f(x)|\ge\varepsilon\}
 \subset \{x:|f_n(x)|\ge\varepsilon/2\}.
\]
右侧因 $f_n\in C_0(X)$ 而紧，左侧又是 $X$ 中的闭集，故它作为右侧的闭子集也紧。于是
$f\in C_0(X)$，并且 $f_n\to f$ 于上确界范数。
\end{proof}

支撑的定义中有一个不能省略的闭包。函数在 $U$ 外为零只推出
$\supp f\subset\overline U$，并不推出 $\supp f\subset U$；
\cref{ex:1-3-1-3}的第一个公式正好展示了差别。另一个边界是完备性：$C_c(X)$ 在上确界范数下通常
不完备。在 $\R$ 上取 $f(x)=(1+|x|)^{-1}$，令 $\chi_n$ 在 $[-n,n]$ 上为 $1$、在
$[-n-1,-n]\cup[n,n+1]$ 上线性降为 $0$、在更外侧为 $0$。则
$f_n=\chi_nf\in C_c(\R)$，且
\[
  \norm{f_n-f}_\infty\le\frac{1}{1+n}\longrightarrow0.
\]
因此 $(f_n)$ 在 $C_c(\R)$ 中是 Cauchy 列，但它唯一可能的上确界范数极限 $f$ 没有紧支撑。
完备的自然空间是 $C_0(X)$。

这里得到的截断函数将用于 Radon 正则性、Riesz--Markov 表示以及 $C_c(X)$ 在 $L^p$ 中的稠密性；
它们共同依赖“紧集内取一、指定开集外取零”的支撑控制。

\begin{exercise}[一点紧化中的收敛]
设 $X$ 非紧且局部紧。证明网 $(x_\alpha)$ 在 $X^*$ 中收敛到 $\infty$，当且仅当对每个紧集
$K\subset X$，它最终落在 $X\setminus K$ 中。
\end{exercise}

\begin{exercise}[重建一点紧化判据]
不引用\cref{proposition:1-3-1-one-point-compactification}的证明，分别从“$X$ 局部紧”和“$X^*$
Hausdorff”出发，构造分离 $x$ 与 $\infty$ 的开集或闭包紧的 $x$-邻域，并证明两个方向互逆。
\end{exercise}

\begin{exercise}[局部紧性的遗传]
证明局部紧 Hausdorff 空间的每个开子空间和闭子空间仍局部紧。再证明 $\Q$ 作为 $\R$ 的子空间
不是局部紧的，由此说明任意子空间不保持局部紧性。
\end{exercise}

\begin{exercise}[有限单位分解的直接核对]
在\cref{proposition:1-3-1-finite-partition-of-unity}的构造中，逐项证明分段定义的 $\varphi_j$ 在
$\partial\{s>0\}$ 连续，并证明其支撑不仅紧，而且包含于指定的 $U_j$。
\end{exercise}

\begin{exercise}[$C_0(X)$ 的完备性]
不引用\cref{proposition:1-3-1-c0-complete}的证明，先证明 $C_b(X)$ 在上确界范数下完备，再把
$C_0(X)$ 识别为 $C_b(X)$ 的闭子空间。
\end{exercise}

局部帽函数和无穷远消失条件给出了可用的函数空间；要在其中识别相对紧集，还必须把逐点有界与邻近点间
的共同振荡控制结合起来。

\subsection{等度连续与 \ArzelaAscoli{} 定理}\label{subsec:1-3-2}

紧性逐点地产生极限并不困难。若对每个 $x\in X$，数集 $\{f(x):f\in\mathcal F\}$ 的闭包紧，
乘积紧性会在 $\prod_{x\in X}\overline{\mathcal F(x)}$ 中给出一个点态聚点。真正的问题是：这个聚点
为何连续，点态收敛又为何能够升级为一致收敛？下面的尖峰把缺少的条件画得很清楚。

Ascoli 的定理是函数空间中对有限维 Heine--Borel 原理的一种早期替代：有限维里“有界”已经很强，
函数空间里还必须控制函数在邻近点之间怎样变化。常见的对角子列证明依靠可分性；下面先用网与乘积紧性
证明一般紧 Hausdorff 版本，再在紧度量空间中给出可计算的有限网格证明。

\begin{example}[移动尖峰]\label{ex:1-3-2-2}
在 $[0,1]$ 上令
\[
  f_n(x)=\max\{1-n|x-1/n|,0\}.
\]
这些函数满足 $0\le f_n\le1$，且 $f_n(1/n)=1$。对 $x=0$，有 $f_n(0)=0$；对固定的 $x>0$，
当 $n>2/x$ 时，
\[
 n|x-1/n|>nx/2>1,
\]
所以 $f_n(x)=0$。因此 $f_n$ 逐点收敛到零函数，但
$\norm{f_n}_\infty=1$，不存在一致收敛到零的子列。事实上它们在 $0$ 处不能共享连续性尺度：
$|1/n-0|\to0$，而 $|f_n(1/n)-f_n(0)|=1$。一致有界没有阻止尖峰越来越窄。
\end{example}

\begin{definition}[等度连续]\label{def:1-3-2-1}
设 $X$ 为拓扑空间，$(Y,d_Y)$ 为度量空间。函数族
$\mathcal F\subset C(X,Y)$ 称为在 $x\in X$ \emph{等度连续}，若对每个
$\varepsilon>0$，存在 $x$ 的邻域 $U$，使
\[
 d_Y(f(y),f(x))<\varepsilon
 \qquad(f\in\mathcal F,\ y\in U).
\]
若它在每个 $x$ 都等度连续，就简称 $\mathcal F$ 等度连续。

当 $X$ 为度量空间且函数为标量值时，记
\[
 \omega_{\mathcal F}(\delta)
 =\sup_{f\in\mathcal F}\sup_{d_X(x,y)\le\delta}|f(x)-f(y)|.
\]
若 $X$ 紧，则等度连续等价于
$\omega_{\mathcal F}(\delta)\to0$（$\delta\downarrow0$）。
\end{definition}

说明最后的等价。共同模量趋于零直接给出各点的等度连续。反过来，给定 $\varepsilon>0$，对每个
$x\in X$ 可取 $r_x>0$，使 $d_X(x,y)<2r_x$ 时所有 $f\in\mathcal F$ 都满足
$|f(y)-f(x)|<\varepsilon/2$。有限多个 $B(x_i,r_{x_i})$ 覆盖 $X$。令
$\delta=\min_i r_{x_i}$。若 $d_X(y,z)\le\delta$，选 $i$ 使 $y\in B(x_i,r_{x_i})$，则
$z\in B(x_i,2r_{x_i})$，从而
\[
 |f(y)-f(z)|
 \le |f(y)-f(x_i)|+|f(x_i)-f(z)|<\varepsilon.
\]

\begin{definition}[点态相对紧]\label{def:1-3-2-2}
若对每个 $x\in X$，集合
$\mathcal F(x)=\{f(x):f\in\mathcal F\}$ 在 $Y$ 中的闭包紧，就称 $\mathcal F$ \emph{点态相对紧}。
标量值函数族点态相对紧，当且仅当它在每一点的取值集合有界。
\end{definition}

\begin{definition}[共同在无穷远消失]\label{def:1-3-2-3}
设 $X$ 局部紧，$\mathcal F\subset C_0(X)$。若对每个 $\varepsilon>0$，存在紧集 $K\subset X$，
使
\[
 |f(x)|<\varepsilon
 \qquad(f\in\mathcal F,\ x\in X\setminus K),
\]
就称 $\mathcal F$ \emph{共同在无穷远消失}。
\end{definition}

若 $X$ 紧而 $Y$ 为度量空间，连续函数 $f,g:X\to Y$ 的距离函数
$x\mapsto d_Y(f(x),g(x))$ 在紧集上有界。因此
\[
  d_\infty(f,g)=\sup_{x\in X}d_Y(f(x),g(x))
\]
是 $C(X,Y)$ 上的度量（$X=\varnothing$ 时仍约定上确界为 $0$）；它诱导的拓扑称为
\emph{一致拓扑}，而 $d_\infty(f_\alpha,f)\to0$ 就是一致收敛。后文紧域上的“一致相对紧”均相对于
这个度量。

\begin{proposition}[等度连续极限升级]\label{proposition:1-3-2-1}
设 $(Y,d_Y)$ 为度量空间，$(f_\alpha)$ 是 $C(X,Y)$ 中的网，并且函数族
$\{f_\alpha\}$ 等度连续。若 $f_\alpha(x)\to f(x)$ 对每个 $x\in X$ 成立，则 $f$ 连续；若 $X$ 紧，
则 $f_\alpha\to f$ 一致收敛。
\end{proposition}

\begin{proof}
固定 $x\in X$ 与 $\varepsilon>0$。等度连续性给出 $x$ 的邻域 $U$，使
\[
 d_Y(f_\alpha(y),f_\alpha(x))<\varepsilon/3
 \qquad(y\in U,\ \alpha).
\]
对固定的 $y\in U$，令 $\alpha$ 趋向有向集的尾部；度量的连续性给出
\[
 d_Y(f(y),f(x))
 =\lim_\alpha d_Y(f_\alpha(y),f_\alpha(x))
 \le\varepsilon/3<\varepsilon.
\]
所以 $f$ 在 $x$ 连续。

再设 $X$ 紧。对每个 $x\in X$，由等度连续性与刚证得的 $f$ 的连续性，存在邻域 $U_x$，使对
$y\in U_x$ 与所有 $\alpha$ 都有
\[
 d_Y(f_\alpha(y),f_\alpha(x))<\varepsilon/3,
 \qquad
 d_Y(f(y),f(x))<\varepsilon/3.
\]
取有限子覆盖 $U_{x_1},\ldots,U_{x_m}$。逐点收敛在这有限多个点上可同时实现：存在 $\alpha_0$，使
$\alpha\ge\alpha_0$ 时
\[
 d_Y(f_\alpha(x_i),f(x_i))<\varepsilon/3
 \qquad(1\le i\le m).
\]
给定 $y\in X$，选 $i$ 使 $y\in U_{x_i}$。三角不等式给出
\[
 d_Y(f_\alpha(y),f(y))
 \le d_Y(f_\alpha(y),f_\alpha(x_i))
   +d_Y(f_\alpha(x_i),f(x_i))
   +d_Y(f(x_i),f(y))
 <\varepsilon.
\]
这个估计对所有 $y$ 同时成立，故收敛一致。
\end{proof}

\begin{theorem}[\ArzelaAscoli{}：紧域]\label{theorem:1-3-2-2}
设 $X$ 为紧 Hausdorff 空间，$(Y,d_Y)$ 为度量空间。若
$\mathcal F\subset C(X,Y)$ 等度连续且点态相对紧，则 $\mathcal F$ 在一致拓扑下相对紧。
特别地，当 $Y$ 紧时，点态相对紧自动成立；标量值情形中，它等价于点态有界。
\end{theorem}

\begin{proof}
对每个 $x\in X$ 令 $K_x=\overline{\mathcal F(x)}\subset Y$。由假设 $K_x$ 紧，故由
Tychonoff 定理，乘积
\[
  E=\prod_{x\in X}K_x
\]
在点态收敛拓扑下紧。把每个 $f\in\mathcal F$ 看作 $E$ 中的点，并令 $H$ 为 $\mathcal F$ 在 $E$
中的闭包；$H$ 是紧集。

先证明 $H$ 仍是等度连续的连续函数族。固定 $h\in H$。由闭包的网刻画，存在
$(f_\alpha)\subset\mathcal F$ 在 $E$ 中收敛到 $h$，也就是逐点收敛到 $h$。由
\cref{proposition:1-3-2-1}，$h$ 连续。再固定 $x\in X$ 与 $\varepsilon>0$。原函数族的等度连续性
给出邻域 $U$，使
\[
 d_Y(f(y),f(x))<\varepsilon/2
 \qquad(f\in\mathcal F,\ y\in U).
\]
对任意 $h\in H$ 取上述逼近网并令其逐点趋于 $h$，便有
\[
 d_Y(h(y),h(x))\le\varepsilon/2<\varepsilon
 \qquad(y\in U).
\]
同一个 $U$ 对所有 $h\in H$ 有效，所以 $H$ 等度连续。

现在考察恒等映射
\[
 \iota:(H,\text{点态拓扑})\longrightarrow(H,\text{一致拓扑}).
\]
若网 $(h_\beta)\subset H$ 在点态拓扑下收敛到 $h\in H$，则刚证得的等度连续性与
\cref{proposition:1-3-2-1}说明它一致收敛到 $h$。由网的连续性判据，$\iota$ 连续。因此点态紧集
$H$ 在一致拓扑下仍紧。

还要确认 $H$ 正是 $\mathcal F$ 的一致闭包。每个 $h\in H$ 都是 $\mathcal F$ 中某网的点态极限，
而该网由\cref{proposition:1-3-2-1}一致收敛到 $h$，故 $H$ 包含于一致闭包。反过来，$H$ 在一致
拓扑下紧，因而在 Hausdorff 的一致拓扑下闭；它包含 $\mathcal F$，所以也包含 $\mathcal F$ 的一致闭包。
两者相等，且这个闭包紧，结论得证。
\end{proof}

乘积紧性只产生坐标极限。上述证明中，等度连续性承担了两项不同工作：先保证极限仍连续，再保证点态拓扑
在闭包上升级为一致拓扑。省略其中任何一步，都会把紧集放在错误的函数空间或错误的拓扑里。

\begin{theorem}[\ArzelaAscoli{}：$C_0$ 版]\label{theorem:1-3-2-3}
设 $X$ 局部紧 Hausdorff。若标量函数族 $\mathcal F\subset C_0(X)$ 等度连续、点态有界，并且共同在
无穷远消失，则 $\mathcal F$ 在上确界范数下相对紧。
\end{theorem}

\begin{proof}
若 $X$ 紧，则 $C_0(X)=C(X)$，结论直接来自紧域版本。以下设 $X$ 非紧。对每个 $f\in\mathcal F$
定义 $\widetilde f\in C(X^*)$，令 $\widetilde f|_X=f$ 且 $\widetilde f(\infty)=0$。
在 $X$ 中各点的等度连续性没有改变；而共同在无穷远消失正好表示：对每个 $\varepsilon>0$，存在
$\infty$ 的邻域 $\{\infty\}\cup(X\setminus K)$，使
\[
 |\widetilde f(x)-\widetilde f(\infty)|<\varepsilon
 \qquad(f\in\mathcal F)
\]
在该邻域成立。因此延拓族在 $X^*$ 上等度连续。它在 $X$ 上点态有界，在 $\infty$ 处恒为零，也就
点态有界。紧域版本说明 $\{\widetilde f:f\in\mathcal F\}$ 在 $C(X^*)$ 中相对紧。延拓保持上确界范数，
而 $\{g\in C(X^*):g(\infty)=0\}$ 是连续评价映射 $g\mapsto g(\infty)$ 的闭核。因此延拓族的闭包
仍在这个闭核中，故 $\mathcal F$ 在 $C_0(X)$ 中相对紧。
\end{proof}

在紧度量空间上，可以把同一个原理写成有限采样算法，不必在证明中显式处理任意乘积。

\begin{proposition}[紧度量空间的有限采样判据]\label{proposition:1-3-2-4}
设 $(X,d)$ 紧，$\mathcal F\subset C(X,\C)$ 点态有界，并且存在函数
$\omega:[0,\infty)\to[0,\infty)$，满足 $\omega(\delta)\to0$（$\delta\downarrow0$）以及
\[
  |f(x)-f(y)|\le\omega(d(x,y))
  \qquad(f\in\mathcal F,\ x,y\in X).
\]
则 $\mathcal F$ 在上确界范数下全有界，因而相对紧。
\end{proposition}

\begin{proof}
若 $X=\varnothing$ 或 $\mathcal F=\varnothing$，结论成立。以下设二者均非空。
给定 $\varepsilon>0$，由 $\omega(t)\to0$ 可选 $\delta>0$，使 $0<t<\delta$ 时都有
$\omega(t)<\varepsilon/3$。紧性给出有限
$\delta$-网 $x_1,\ldots,x_m$。采样映射
\[
  S(f)=(f(x_1),\ldots,f(x_m))\in\C^m
\]
的像有界，因为只涉及有限多个点，而函数族在每一点有界。有限维 Heine--Borel 定理说明
$S(\mathcal F)$ 全有界。以下在 $\C^m$ 上使用最大值范数。先取它的有限 $\varepsilon/6$-网；
对每个与 $S(\mathcal F)$ 相交的
网球，从交集中选一个采样向量，并取一个原像函数 $g_k\in\mathcal F$。这里只有有限多个网球。
三角不等式说明每个 $f\in\mathcal F$ 都有某个 $g_k$ 满足
\[
  \max_{1\le i\le m}|f(x_i)-g_k(x_i)|<\varepsilon/3.
\]
对任意 $x\in X$，取 $i$ 使 $d(x,x_i)<\delta$。若 $x=x_i$，下式首尾两个模量误差均为零；
否则二者都小于 $\varepsilon/3$。于是
\[
 |f(x)-g_k(x)|
 \le |f(x)-f(x_i)|+|f(x_i)-g_k(x_i)|+|g_k(x_i)-g_k(x)|
 <\varepsilon.
\]
所以 $g_1,\ldots,g_N$ 是 $\mathcal F$ 的有限 $\varepsilon$-网。由 Analysis II 中完备度量空间的
全有界判据，$C(X,\C)$ 的完备性把全有界性升级为相对紧性。
\end{proof}

\begin{figure}[htbp]
\centering
\begin{tikzpicture}[node distance=12mm and 13mm]
  \node[book strong node] (grid) {有限 $\delta$-网\\$x_1,\ldots,x_m$};
  \node[book node,right=of grid] (sample) {采样向量\\$S(f)\in\C^m$};
  \node[book node,right=of sample] (reps) {有限代表\\$g_1,\ldots,g_N$};
  \node[book warm node,below=of reps] (uniform) {$\norm{f-g_k}_\infty<\varepsilon$};
  \node[book node,left=22mm of uniform] (modulus) {共同模量\\$\omega(t)<\varepsilon/3\quad(0<t<\delta)$};
  \draw[book arrow] (grid) -- node[above,book label]{观测} (sample);
  \draw[book arrow] (sample) -- node[above,book label]{$\C^m$ 全有界} (reps);
  \draw[book arrow] (reps) -- (uniform);
  \draw[book accent arrow] (modulus) -- node[below,book label]{误差传播} (uniform);
\end{tikzpicture}
\caption{有限采样控制采样向量，共同模量把有限坐标上的误差传播到整个紧集。}
\label{fig:1-3-2-finite-sampling}
\end{figure}

\begin{example}[Lipschitz 球]\label{ex:1-3-2-1}
设 $(X,d)$ 紧，固定 $L,M\ge0$，并令
\[
 \mathcal L_{L,M}
 =\{f\in C(X,\R):|f(x)|\le M,
   \ |f(x)-f(y)|\le Ld(x,y)\text{ 对所有 }x,y\}.
\]
它点态有界，并有共同模量 $\omega(\delta)=L\delta$。由
\cref{proposition:1-3-2-4}，$\mathcal L_{L,M}$ 在 $C(X)$ 中相对紧。它还关于一致收敛闭：若
$f_n\to f$ 一致，则逐点取极限得到 $|f|\le M$ 与
$|f(x)-f(y)|\le Ld(x,y)$。所以 $\mathcal L_{L,M}$ 本身是紧集。
\end{example}

非紧域上还需控制函数向无穷远移动。取
$\phi(x)=\max\{1-|x|,0\}$ 并令 $f_n(x)=\phi(x-3n)$。函数族
$\{f_n\}\subset C_0(\R)$ 共同 $1$-Lipschitz，且在每一点只有至多一个函数非零，因此点态有界；但它不
共同在无穷远消失。不同函数的支撑不交且峰值均为 $1$，故
$\norm{f_n-f_m}_\infty=1$（$n\ne m$）。它没有 Cauchy 子列，从而不相对紧。

点态有界本身也不等于一致有界。令
\[
  c_n=\frac12\left(\frac1{n+1}+\frac1n\right),\qquad
  r_n=\frac12\left(\frac1n-\frac1{n+1}\right),
\]
并在 $[0,1]$ 上定义
\[
  h_n(x)=n\max\left\{1-\frac{|x-c_n|}{r_n},0\right\}.
\]
则 $h_n$ 连续，非零集包含于互不相交的区间 $(1/(n+1),1/n)$，且
\[
  \norm{h_n}_\infty=h_n(c_n)=n.
\]
固定 $x$ 后，至多一个 $h_n(x)$ 非零，所以
$\sup_n|h_n(x)|<\infty$；但上确界范数没有共同界。等度连续性恰好排除了这种越来越陡的峰。

\begin{remark}[积分原函数的共同模量]\label{ex:1-3-2-3}
设 $p>1$、$q=p/(p-1)$，并设函数 $g_n\in L^p([0,1])$ 满足
$\norm{g_n}_{L^p}\le M$。令
\[
  F_n(x)=\int_0^xg_n(t)\,\dd t
\]
。对任意 $x,y\in[0,1]$，记
$I_{x,y}=[\min\{x,y\},\max\{x,y\}]$。\Holder{} 不等式给出
\[
 |F_n(y)-F_n(x)|
 \le \norm{g_n}_{L^p}\norm{\1_{I_{x,y}}}_{L^q}
 \le M|x-y|^{1/q}.
\]
这给出共同 \Holder{} 模量；取 $y=0$ 还得到
$|F_n(x)|\le Mx^{1/q}\le M$，故函数族点态有界。\ArzelaAscoli{} 定理因而给出一致收敛子列。
本评注只说明该紧性定理在积分族上的一种典型应用；$L^p$ 积分与 \Holder{} 不等式见第~3~章。
\end{remark}

同样的紧性结构以后会出现在常微分方程近似解、随机过程的路径以及把有界集合送入相对紧函数族的算子中。
经验过程与 Kolmogorov 连续性定理会再次使用这种机制；Helly 选择原理则依靠单调性或有界变差而不是
等度连续性。上述对象各自还需要新的定义。\ArzelaAscoli{} 只提供函数空间中的紧性判据，并不替代那些
后续论证。

\begin{exercise}[闭包保持等度连续]
设 $X$ 为拓扑空间、$Y$ 为度量空间。证明等度连续族在点态闭包中的所有连续函数仍组成等度连续族；
再说明在紧 $X$ 上，点态闭包与一致闭包在 \ArzelaAscoli{} 的假设下为何相同。
\end{exercise}

\begin{exercise}[一点值控制]
设 $(X,d)$ 紧，函数族共同 $L$-Lipschitz。证明只要在某一点 $x_0$ 的取值集合有界，函数族就在
$X$ 上一致有界，并由此在 $C(X)$ 中相对紧。
\end{exercise}

\begin{exercise}[导数控制的函数族]
证明
\[
 \{f\in C^1([0,1]):|f(0)|\le M,\ \norm{f'}_\infty\le M\}
\]
在 $C([0,1])$ 中相对紧。若删去 $|f(0)|\le M$，说明常数函数为何给出反例。
\end{exercise}

\begin{exercise}[非紧域的缺失条件]
逐项核对平移函数 $f_n(x)=\max\{1-|x-3n|,0\}$ 的等度连续性、点态有界性与
$C_0(\R)$ 成员资格，并直接证明它们不共同在无穷远消失且没有一致收敛子列。
\end{exercise}

有限采样定理告诉我们怎样压缩一个已知函数族。反过来，从少量基本观测函数出发生成整个连续函数空间，
需要把点分离与代数运算转化为一致逼近。

\subsection{Stone--Weierstrass：分离观测生成函数}\label{subsec:1-3-3}

\ArzelaAscoli{} 从有限采样控制一个已经给定的函数族。现在把问题反过来：给定一族基本观测函数，允许作
有限次加法、乘法与极限，它们能否生成所有连续函数？区间上的多项式暗示答案与坐标函数有关，但复数情形
马上提醒我们，单纯“能分离点”还少一个条件。

经典 Weierstrass 定理只处理区间多项式。Stone 的抽象化指出，多项式的具体公式不是核心：含常数、
分离点以及复情形中的共轭封闭，才是把局部插值拼成全局逼近的结构。下面的证明会把这一抽象重新拆回
绝对值、最大值、最小值与两次有限覆盖。

\begin{example}[缺少共轭的失败]\label{ex:1-3-3-3}
令 $\overline{\mathbb D}=\{z\in\C:|z|\le1\}$，并把圆盘代数
$A(\overline{\mathbb D})$ 定义为复多项式在 $C(\overline{\mathbb D},\C)$ 中的一致闭包。它含常数，
坐标函数 $z\mapsto z$ 已经分离闭圆盘中的点。然而 $\overline z$ 不属于这个闭包。

证明只用实变量 Riemann 积分。若多项式 $p_n$ 在闭圆盘上一致逼近 $\overline z$，则在单位圆上
$p_n(e^{it})e^{it}$ 一致趋于 $1$。把复值 Riemann 积分定义为实部与虚部的两个实积分，有
\[
 \left|\int_0^{2\pi}p_n(e^{it})e^{it}\,\dd t-2\pi\right|
 \le 2\pi\sup_{0\le t\le2\pi}
       |p_n(e^{it})-e^{-it}|\longrightarrow0.
\]
另一方面，若 $p(z)=\sum_{k=0}^Na_kz^k$，则
\[
 \int_0^{2\pi}p(e^{it})e^{it}\,\dd t
 =\sum_{k=0}^Na_k\int_0^{2\pi}e^{i(k+1)t}\,\dd t=0,
\]
因为对每个非零整数 $m$，$\cos(mt)$ 与 $\sin(mt)$ 在 $[0,2\pi]$ 上的积分都为零。于是左侧对每个
$n$ 都等于零，不可能趋于 $2\pi$。这个反例没有使用 Cauchy 定理、围道积分或解析函数理论。
\end{example}

\begin{definition}[分离点与处处不消失]\label{def:1-3-3-1}
函数族 $S\subset C(X,\C)$ 称为\emph{分离点}，若对任意 $x\ne y$，存在 $f\in S$ 使
$f(x)\ne f(y)$。若对每个 $x\in X$，存在 $f\in S$ 使 $f(x)\ne0$，就称 $S$ \emph{处处不消失}。
\end{definition}

\begin{definition}[函数代数与 $*$-子代数]\label{def:1-3-3-2}
设 $\mathbb F\in\{\R,\C\}$。函数代数 $A\subset C(X,\mathbb F)$ 是对加法、标量乘法与逐点乘法
封闭的线性子空间；除非另行说明，不要求它含常数。若复函数代数还满足
\[
  f\in A\quad\Longrightarrow\quad\overline f\in A,
\]
就称它为 $*$-子代数。这里 $\overline f(x)=\overline{f(x)}$。实函数代数自动满足对应条件。
\end{definition}

\begin{definition}[一致闭包与函数格]\label{def:1-3-3-3}
函数族 $A\subset C_b(X)$ 关于上确界范数的闭包称为它的\emph{一致闭包}，记作 $\overline A$；
当 $X$ 紧时 $C_b(X)=C(X)$。实函数
线性空间 $L$ 若对逐点的 $\max$ 与 $\min$ 封闭，就称为\emph{函数格}。
\end{definition}

把代数变成函数格是证明的关键。看似需要先知道多项式能逼近绝对值，但若直接调用经典
Weierstrass 定理，就会用待证结论的特例证明一般结论。下面从实幂级数独立构造所需多项式。

\begin{lemma}[绝对值进入闭包]\label{lemma:1-3-3-1}
设 $X$ 紧 Hausdorff，$A\subset C(X,\R)$ 是含常数的函数代数。若 $f\in\overline A$，则
$|f|\in\overline A$。因而 $\overline A$ 是函数格。
\end{lemma}

\begin{proof}
先记 $B=\overline A$。它仍是含常数的函数代数。加法与标量乘法的封闭性来自范数连续性；对于乘法，
若 $f_n\to f$、$g_n\to g$ 一致，则 $(g_n)$ 一致有界，并且
\[
 \norm{f_ng_n-fg}_\infty
 \le\norm{f_n-f}_\infty\norm{g_n}_\infty
   +\norm f_\infty\norm{g_n-g}_\infty\longrightarrow0.
\]

令 $M=\norm f_\infty$。$M=0$ 时结论成立。设 $M>0$。实变量二项级数给出
\begin{equation}\label{eq:1-3-binomial-sqrt}
 \sqrt{1-s}=1-\sum_{n=1}^\infty c_ns^n,
 \qquad 0\le s<1,
\end{equation}
其中
\[
 c_n=\frac{1}{2n-1}\frac{\binom{2n}{n}}{4^n}>0.
\]
还需确认端点 $s=1$ 不破坏一致逼近。对 $0<r<1$，
$\sum_{n\ge1}c_nr^n=1-\sqrt{1-r}\le1$。任意有限部分和令 $r\uparrow1$ 后不超过 $1$，故
$\sum_nc_n\le1$；反过来
$\sum_nc_n\ge\sum_nc_nr^n=1-\sqrt{1-r}$，再令 $r\uparrow1$ 得
$\sum_nc_n\ge1$。所以 $\sum_nc_n=1$。由正项级数的尾和估计，
\eqref{eq:1-3-binomial-sqrt}在 $[0,1]$ 上一致收敛。

定义实多项式
\[
 q_N(t)=M\left[1-\sum_{n=1}^Nc_n
             \left(1-\frac{t^2}{M^2}\right)^n\right].
\]
当 $|t|\le M$ 时，取 $s=1-t^2/M^2\in[0,1]$，由上面的均匀尾和有
\[
 \sup_{|t|\le M}|q_N(t)-|t||
 \le M\sum_{n>N}c_n\longrightarrow0.
\]
因为 $B$ 是含常数代数，$q_N(f)\in B$；一致取极限得到 $|f|\in B$。最后用恒等式
\[
 \max(f,g)=\frac{f+g+|f-g|}{2},
 \qquad
 \min(f,g)=\frac{f+g-|f-g|}{2}
\]
可知 $B$ 对有限的逐点最大值和最小值封闭。
\end{proof}

\begin{theorem}[Stone--Weierstrass：实情形]\label{theorem:1-3-3-2}
设 $X$ 紧 Hausdorff。若 $A\subset C(X,\R)$ 是含常数且分离点的函数代数，则 $A$ 在
$C(X,\R)$ 中一致稠密。
\end{theorem}

\begin{proof}
若 $X=\varnothing$，则 $C(X,\R)$ 只有零函数，结论成立。以下设 $X\ne\varnothing$。
记 $B=\overline A$。由\cref{lemma:1-3-3-1}，$B$ 是函数格。固定目标函数
$f\in C(X,\R)$ 与 $\varepsilon>0$。

对 $x,y\in X$ 构造 $h_{x,y}\in A$，使
\[
 h_{x,y}(x)=f(x),\qquad h_{x,y}(y)=f(y).
\]
当 $x=y$ 时取常数 $f(x)$。当 $x\ne y$ 时，由分离点可取 $a\in A$ 使 $a(x)\ne a(y)$，并置
\[
 h_{x,y}
 =f(x)+\frac{f(y)-f(x)}{a(y)-a(x)}\,[a-a(x)].
\]
因为 $A$ 含常数，$h_{x,y}\in A$，且两个插值等式成立。

先固定 $x$。对每个 $y\in X$，连续函数 $h_{x,y}-f$ 在 $y$ 处为零，所以开集
\[
 V_{x,y}=\{z:h_{x,y}(z)>f(z)-\varepsilon\}
\]
包含 $y$。这些开集覆盖 $X$，由紧性选出 $y_1,\ldots,y_m$ 使
$V_{x,y_1},\ldots,V_{x,y_m}$ 覆盖 $X$。令
\[
 g_x=\max_{1\le j\le m}h_{x,y_j}\in B.
\]
对每个 $z\in X$，至少一个 $V_{x,y_j}$ 包含 $z$，故
$g_x(z)>f(z)-\varepsilon$。另一方面，所有 $h_{x,y_j}(x)$ 都等于 $f(x)$，所以 $g_x(x)=f(x)$。
于是
\[
 U_x=\{z:g_x(z)<f(z)+\varepsilon\}
\]
是包含 $x$ 的开集。

当 $x$ 变化时，$U_x$ 覆盖 $X$。再由紧性选取 $x_1,\ldots,x_r$ 使这些开集覆盖 $X$，并令
\[
  g=\min_{1\le i\le r}g_{x_i}\in B.
\]
每个 $g_{x_i}$ 在全空间都大于 $f-\varepsilon$，故
$g>f-\varepsilon$。给定 $z\in X$，选 $i$ 使 $z\in U_{x_i}$，则
$g(z)\le g_{x_i}(z)<f(z)+\varepsilon$。因此
\[
  |g(z)-f(z)|<\varepsilon\qquad(z\in X).
\]
这说明 $f$ 属于 $B$。目标函数任意，故 $B=C(X,\R)$。
\end{proof}

证明中的两次有限覆盖承担不同方向的不等式。第一次对固定 $x$ 取有限最大值，使函数在全空间位于
$f-\varepsilon$ 上方；第二次取有限最小值，使每一点至少受到一个 $f+\varepsilon$ 上界的控制。
若交换最大值与最小值，两个方向不会自动保留。

\begin{theorem}[Stone--Weierstrass：复情形]\label{theorem:1-3-3-3}
设 $X$ 紧 Hausdorff。若 $A\subset C(X,\C)$ 是含常数、分离点的 $*$-子代数，则 $A$ 在
$C(X,\C)$ 中一致稠密。
\end{theorem}

\begin{proof}
令
\[
 A_{\R}=\{a\in A:a=\overline a\}.
\]
它是 $C(X,\R)$ 中含常数的实函数代数。若 $x\ne y$，取 $a\in A$ 使 $a(x)\ne a(y)$。两个复数
$a(x),a(y)$ 的实部或虚部至少有一个不同，而
\[
 \Re a=\frac{a+\overline a}{2}\in A_{\R},
 \qquad
 \Im a=\frac{a-\overline a}{2i}\in A_{\R}.
\]
所以 $A_{\R}$ 分离点。由实情形，$A_{\R}$ 在 $C(X,\R)$ 中稠密。

给定 $f=u+iv\in C(X,\C)$ 与 $\varepsilon>0$，可取 $a,b\in A_{\R}$ 使
$\norm{u-a}_\infty<\varepsilon/2$、$\norm{v-b}_\infty<\varepsilon/2$。于是 $a+ib\in A$，且
\[
 \norm{f-(a+ib)}_\infty
 \le\norm{u-a}_\infty+\norm{v-b}_\infty<\varepsilon.
\]
故 $A$ 一致稠密。
\end{proof}

\begin{figure}[htbp]
\centering
\begin{tikzpicture}[node distance=11mm and 13mm]
  \node[book strong node] (sep) {分离 $x,y$};
  \node[book node,right=of sep] (interp) {两点插值\\$h_{x,y}$};
  \node[book node,right=of interp] (max) {有限 $\max$\\固定 $x$};
  \node[book node,below=of max] (min) {有限 $\min$\\覆盖 $X$};
  \node[book warm node,left=of min] (approx) {$\norm{g-f}_\infty<\varepsilon$};
  \draw[book arrow] (sep) -- (interp);
  \draw[book arrow] (interp) -- node[above,book label]{格运算} (max);
  \draw[book arrow] (max) -- node[right,book label]{紧性} (min);
  \draw[book accent arrow] (min) -- (approx);
\end{tikzpicture}
\caption{Stone--Weierstrass 证明的两次拼接。最大值与最小值分别保留一侧误差。}
\label{fig:1-3-3-stone-weierstrass}
\end{figure}

\begin{example}[区间多项式]\label{ex:1-3-3-1}
实多项式在 $[a,b]$ 上构成含常数的函数代数。若 $x\ne y$，坐标多项式 $p(t)=t$ 满足
$p(x)\ne p(y)$，所以它分离点。由\cref{theorem:1-3-3-2}，对每个
$f\in C([a,b],\R)$ 与 $\varepsilon>0$，存在实多项式 $p$ 使
$\norm{f-p}_\infty<\varepsilon$。经典 Weierstrass 逼近定理由此成为一般定理的直接实例。
\end{example}

\begin{example}[三角多项式]\label{ex:1-3-3-2}
在 $\mathbb T=\{z\in\C:|z|=1\}$ 上，坐标函数 $\zeta(z)=z$ 与其共轭满足
$\overline\zeta(z)=z^{-1}$。它们生成的 $*$-代数恰为有限 Laurent 和
\[
  \sum_{k=-N}^Na_kz^k,
\]
也就是三角多项式。该代数含常数，且 $\zeta$ 本身分离 $\mathbb T$ 的点。由
\cref{theorem:1-3-3-3}，三角多项式在 $C(\mathbb T,\C)$ 中一致稠密。这里的结论仍然只涉及
连续函数与代数运算；Fourier 系数和复积分尚未出现。
\end{example}

圆盘反例现在也可准确定位：复多项式闭包含常数并分离点，却不对共轭封闭；正是缺少的 $*$ 条件使
$\overline z$ 无法进入闭包。

\begin{theorem}[Stone--Weierstrass：局部紧版本]\label{theorem:1-3-3-4}
设 $X$ 局部紧 Hausdorff，$A\subset C_0(X,\C)$ 是分离点、处处不消失的 $*$-子代数。则 $A$ 在
$C_0(X,\C)$ 中一致稠密。
\end{theorem}

\begin{proof}
先设 $X$ 非紧。把每个 $a\in A$ 以 $a(\infty)=0$ 延拓到 $X^*$，所得代数记为 $\widetilde A$，并令
\[
  A^+=\widetilde A+\C\mathbf1\subset C(X^*,\C).
\]
这是含常数的 $*$-子代数。$A$ 分离 $X$ 内的两点；若 $x\in X$，处处不消失给出 $a\in A$ 使
$a(x)\ne0=a(\infty)$，所以 $A^+$ 也分离 $x$ 与 $\infty$。紧空间复版本说明 $A^+$ 在
$C(X^*)$ 中稠密。

给定 $f\in C_0(X)$，把它延拓为 $\widetilde f\in C(X^*)$，其中 $\widetilde f(\infty)=0$。给定
$\varepsilon>0$，取 $a\in A$ 与 $c\in\C$，使
\[
 \norm{\widetilde f-(\widetilde a+c\mathbf1)}_\infty<\varepsilon/2.
\]
在 $\infty$ 处代入得到 $|c|<\varepsilon/2$，故
\[
 \norm{f-a}_\infty
 \le\norm{\widetilde f-(\widetilde a+c\mathbf1)}_\infty+|c|
 <\varepsilon.
\]

若 $X$ 本身紧，就令 $Y=X\sqcup\{\infty\}$，并把新点规定为孤立点。$Y$ 是紧 Hausdorff 空间，
但这里不把它称为 $X$ 的一点紧化，因为 $X$ 在 $Y$ 中不稠密。把 $A$ 中函数在新点取零，并把所得代数
记为 $\widetilde A$。代数 $\widetilde A+\C\mathbf1$ 用 $A$ 分离 $X$ 中的两点；给定 $x\in X$，
处处不消失性给出 $a\in A$ 使 $a(x)\ne0=\widetilde a(\infty)$，所以它也分离 $x$ 与 $\infty$。
紧空间复版本表明该含常数 $*$-代数在 $C(Y)$ 中稠密。把 $f\in C(X)=C_0(X)$ 在 $\infty$ 处延拓为
零，并取
\(\norm{\widetilde f-(\widetilde a+c\mathbf1)}_\infty<\varepsilon/2\)。在 $\infty$ 处代入得
$|c|<\varepsilon/2$，限制回 $X$ 后便有
$\norm{f-a}_\infty<\varepsilon$，结论成立。
\end{proof}

局部紧版本不要求 $A$ 含常数。在非紧空间中，非零常数通常不属于 $C_0(X)$；处处不消失取代了含常数
所承担的一部分作用，因为它恰好把空间中的每个点与 $\infty$ 分开。分离点也只是全局的观测条件，
与导数是否非零、是否提供局部坐标没有等价关系。

\begin{corollary}[紧度量空间上的可数稠密测试族]\label{corollary:1-3-3-separable}
若 $(X,d)$ 是紧度量空间，则 $C(X,\R)$ 与 $C(X,\C)$ 都可分。
\end{corollary}

\begin{proof}
对每个 $n$ 取有限 $1/n$-网，所有这些网的并给出可数稠密集 $D\subset X$；按后文的记号，这个
逐个选取有限网的步骤使用 $\textsf{CC}$。对 $q\in D$ 定义
$h_q(x)=d(x,q)$。这族函数分离点：若 $x\ne y$，令 $\rho=d(x,y)>0$，取
$q\in D$ 使 $d(x,q)<\rho/3$，则
\[
 |h_q(y)-h_q(x)|
 \ge d(x,y)-2d(x,q)>\rho/3.
\]
由 $1$ 与所有 $h_q$ 生成的实函数代数依
\cref{theorem:1-3-3-2}在 $C(X,\R)$ 中稠密。只允许有理系数、有限多个生成元的多项式仍构成可数集。
具体地，固定
\[
  P=\sum_{\alpha\in F}c_\alpha
       h_{q_1}^{\alpha_1}\cdots h_{q_m}^{\alpha_m},
\]
其中 $F\subset\N^m$ 有限，并给定 $\varepsilon>0$。令
$M_\alpha=\prod_{j=1}^m\norm{h_{q_j}}_\infty^{\alpha_j}$（约定 $0^0=1$），并选
$r_\alpha\in\Q$ 使
\[
  \sum_{\alpha\in F}|c_\alpha-r_\alpha|M_\alpha<\varepsilon.
\]
则相应有理系数多项式 $Q$ 满足
\[
  \norm{P-Q}_\infty
  \le\sum_{\alpha\in F}|c_\alpha-r_\alpha|M_\alpha<\varepsilon.
\]
故这个可数子集已经在 $C(X,\R)$ 中稠密。复情形改用有理实部与有理虚部的系数即可。
\end{proof}

以下只作后文术语预告；Radon 测度与积分将在测度论章节定义，本章不以这段话为任何证明的前提。
对两个有限 Radon 测度 $\mu,\nu$，若它们在一个含常数、分离点的 $*$-代数 $A$ 上给出相同积分，
则 Stone--Weierstrass 定理允许对任意连续函数 $f$ 取 $g_n\in A$ 使
$\norm{g_n-f}_\infty\to0$。有限性给出显式估计
\[
 \left|\int f\,\mathrm d\mu-\int f\,\mathrm d\nu\right|
 \le \mu(X)\norm{f-g_n}_\infty
     +\nu(X)\norm{f-g_n}_\infty
     +\left|\int g_n\,\mathrm d\mu-\int g_n\,\mathrm d\nu\right|.
\]
最后一项为零，令 $n\to\infty$ 即把积分相等性推广到全部连续函数。这一识别方法将在相应概念建立后
用于矩问题与 Fourier 分析。

\begin{exercise}[稠密性的必要条件]
设 $X$ 局部紧 Hausdorff。证明若 $A\subset C_0(X)$ 一致稠密，则 $A$ 分离点且处处不消失。
反过来，指出还必须给 $A$ 增加哪些代数封闭条件，才能应用局部紧 Stone--Weierstrass 定理。
\end{exercise}

\begin{exercise}[三角多项式]
不省略代数验证，证明所有有限和 $\sum_{k=-N}^Na_kz^k$ 在 $\mathbb T$ 上构成含常数、分离点的
$*$-子代数，并由复 Stone--Weierstrass 定理推出其一致稠密性。
\end{exercise}

\begin{exercise}[多变量多项式]
设 $K=\prod_{j=1}^n[a_j,b_j]\subset\R^n$。证明实多变量多项式限制到 $K$ 后构成含常数且分离点的
函数代数，从而在 $C(K,\R)$ 中一致稠密。
\end{exercise}

\begin{exercise}[可数分离族与可度量化]
证明紧 Hausdorff 空间 $X$ 可度量，当且仅当存在可数族 $\{f_n\}\subset C(X,\R)$ 分离点。
提示：一个方向使用距离函数；另一个方向把 $X$ 映入可数乘积
$\prod_n[-\norm{f_n}_\infty,\norm{f_n}_\infty]$，并证明连续单射在紧空间到 Hausdorff 空间之间
给出到其像的同胚。
\end{exercise}

本章还需回答一个隐藏在上述证明里的问题。乘积紧性使用怎样的选择原则？何时可用对角序列代替网？
完备性、可分性与可数拓扑基将决定这些替代何时成立。下一节就整理这些可数性边界。

\input{chapters/ch01/sec-01-04}
